%%% Hlavní soubor. Zde se definují základní parametry a odkazuje se na ostatní části. %%%

%% Verze pro jednostranný tisk:
% Okraje: levý 40mm, pravý 25mm, horní a dolní 25mm
% (ale pozor, LaTeX si sám přidává 1in)
\documentclass[12pt,a4paper]{report}
\setlength\textwidth{145mm}
\setlength\textheight{247mm}
\setlength\oddsidemargin{15mm}
\setlength\evensidemargin{15mm}
\setlength\topmargin{0mm}
\setlength\headsep{0mm}
\setlength\headheight{0mm}
% \openright zařídí, aby následující text začínal na pravé straně knihy
\let\openright=\clearpage

%% Pokud tiskneme oboustranně:
% \documentclass[12pt,a4paper,twoside,openright]{report}
% \setlength\textwidth{145mm}
% \setlength\textheight{247mm}
% \setlength\oddsidemargin{14.2mm}
% \setlength\evensidemargin{0mm}
% \setlength\topmargin{0mm}
% \setlength\headsep{0mm}
% \setlength\headheight{0mm}
% \let\openright=\cleardoublepage

%% Vytváříme PDF/A-2u
\usepackage[a-2u]{pdfx}

%% Přepneme na českou sazbu a fonty Latin Modern
\usepackage[czech]{babel}
\usepackage{lmodern}
\usepackage[T1]{fontenc}
\usepackage{textcomp}

%% Použité kódování znaků: obvykle latin2, cp1250 nebo utf8:
\usepackage[utf8]{inputenc}

%%% Další užitečné balíčky (jsou součástí běžných distribucí LaTeXu)
\usepackage{amsmath}        % rozšíření pro sazbu matematiky
\usepackage{amsfonts}       % matematické fonty
\usepackage{amsthm}         % sazba vět, definic apod.
\usepackage{bbding}         % balíček s nejrůznějšími symboly
			    % (čtverečky, hvězdičky, tužtičky, nůžtičky, ...)
\usepackage{bm}             % tučné symboly (příkaz \bm)
\usepackage{graphicx}       % vkládání obrázků
\graphicspath{{../img/}}
\usepackage{fancyvrb}       % vylepšené prostředí pro strojové písmo
\usepackage{indentfirst}    % zavede odsazení 1. odstavce kapitoly
\usepackage{natbib}         % zajištuje možnost odkazovat na literaturu
			    % stylem AUTOR (ROK), resp. AUTOR [ČÍSLO]
\usepackage[nottoc]{tocbibind} % zajistí přidání seznamu literatury,
                            % obrázků a tabulek do obsahu
\usepackage{icomma}         % inteligetní čárka v matematickém módu
\usepackage{dcolumn}        % lepší zarovnání sloupců v tabulkách
\usepackage{booktabs}       % lepší vodorovné linky v tabulkách
\usepackage{paralist}       % lepší enumerate a itemize
\usepackage{xcolor}         % barevná sazba

\usepackage{xr}
\usepackage{microtype}
\usepackage{subfig}
\usepackage{gensymb}
\renewcommand\_{\textunderscore\allowbreak}
\def\::{\discretionary{::}{}{::}}

%%% Údaje o práci

% Název práce v jazyce práce (přesně podle zadání)
\def\NazevPrace{Analýza videozáznamu skoku o~tyči}

% Název práce v angličtině
\def\NazevPraceEN{Analysis of pole-vault video recording}

% Jméno autora
\def\AutorPrace{Matěj Ščerba}

% Rok odevzdání
\def\RokOdevzdani{2021}

% Název katedry nebo ústavu, kde byla práce oficiálně zadána
% (dle Organizační struktury MFF UK, případně plný název pracoviště mimo MFF)
\def\Katedra{Katedra softwarového inženýrství}
\def\KatedraEN{Department of Software Engineering}

% Jedná se o katedru (department) nebo o ústav (institute)?
\def\TypPracoviste{Katedra}
\def\TypPracovisteEN{Department}

% Vedoucí práce: Jméno a příjmení s~tituly
\def\Vedouci{Ing. Michal Bartoš, Ph.D.}

% Pracoviště vedoucího (opět dle Organizační struktury MFF)
\def\KatedraVedouciho{Katedra softwarového inženýrství}
\def\KatedraVedoucihoEN{Department of Software Engineering}

% Studijní program a obor
\def\StudijniProgram{Informatika (B1801)}
\def\StudijniObor{IOI (1801R008)}

% Nepovinné poděkování (vedoucímu práce, konzultantovi, tomu, kdo
% zapůjčil software, literaturu apod.)
\def\Podekovani{%
Poděkování.
}

% Abstrakt (doporučený rozsah cca 80-200 slov; nejedná se o zadání práce)
\def\Abstrakt{%
Skok o~tyči je jedna z~technicky náročnějších lehkoatletických disciplín, proto se analýzou skoků na většině mezinárodních závodů zabývají vědci zkoumající pohyb lidského těla. Pro trenéry a~závodníky by bylo vhodné, aby měli k~dispozici podobné prostředky i~při trénincích, aniž by museli investovat do specializované techniky. Má práce tuto problematiku řeší. S~použitím počítače a~kamery lze z~jediného videozáznamu získat model atleta a~zároveň analyzovat parametry jím provedeného skoku. Výsledné parametry se posléze zanesou do grafů a~jejich hodnoty lze zkoumat i v~různých fázích skoku za použití prohlížeče videozáznamu.
}
\def\AbstraktEN{%
Abstract.
}

% 3 až 5 klíčových slov (doporučeno), každé uzavřeno ve složených závorkách
\def\KlicovaSlova{%
{skok o~tyči} {detekce pozice těla} {videozáznam}
}
\def\KlicovaSlovaEN{%
{pole vault} {human pose detection} {video recording}
}

%% Balíček hyperref, kterým jdou vyrábět klikací odkazy v PDF,
%% ale hlavně ho používáme k uložení metadat do PDF (včetně obsahu).
%% Většinu nastavítek přednastaví balíček pdfx.
\hypersetup{unicode}
\hypersetup{breaklinks=true}

%% Definice různých užitečných maker (viz popis uvnitř souboru)
%%% Tento soubor obsahuje definice různých užitečných maker a prostředí %%%
%%% Další makra připisujte sem, ať nepřekáží v ostatních souborech.     %%%

%%% Drobné úpravy stylu

% Tato makra přesvědčují mírně ošklivým trikem LaTeX, aby hlavičky kapitol
% sázel příčetněji a nevynechával nad nimi spoustu místa. Směle ignorujte.
\makeatletter
\def\@makechapterhead#1{
  {\parindent \z@ \raggedright \normalfont
   \Huge\bfseries \thechapter. #1
   \par\nobreak
   \vskip 20\p@
}}
\def\@makeschapterhead#1{
  {\parindent \z@ \raggedright \normalfont
   \Huge\bfseries #1
   \par\nobreak
   \vskip 20\p@
}}
\makeatother

% Toto makro definuje kapitolu, která není očíslovaná, ale je uvedena v obsahu.
\def\chapwithtoc#1{
\chapter*{#1}
\addcontentsline{toc}{chapter}{#1}
}

% Trochu volnější nastavení dělení slov, než je default.
\lefthyphenmin=2
\righthyphenmin=2

% Zapne černé "slimáky" na koncích řádků, které přetekly, abychom si
% jich lépe všimli.
\overfullrule=1mm

%%% Makra pro definice, věty, tvrzení, příklady, ... (vyžaduje baliček amsthm)

\theoremstyle{plain}
\newtheorem{veta}{Věta}
\newtheorem{lemma}[veta]{Lemma}
\newtheorem{tvrz}[veta]{Tvrzení}

\theoremstyle{plain}
\newtheorem{definice}{Definice}

\theoremstyle{remark}
\newtheorem*{dusl}{Důsledek}
\newtheorem*{pozn}{Poznámka}
\newtheorem*{prikl}{Příklad}

%%% Prostředí pro důkazy

\newenvironment{dukaz}{
  \par\medskip\noindent
  \textit{Důkaz}.
}{
\newline
\rightline{$\qedsymbol$}
}

%%% Prostředí pro sazbu kódu, případně vstupu/výstupu počítačových
%%% programů. (Vyžaduje balíček fancyvrb -- fancy verbatim.)

\DefineVerbatimEnvironment{code}{Verbatim}{fontsize=\small, frame=single}

%%% Prostor reálných, resp. přirozených čísel
\newcommand{\R}{\mathbb{R}}
\newcommand{\N}{\mathbb{N}}

%%% Užitečné operátory pro statistiku a pravděpodobnost
\DeclareMathOperator{\pr}{\textsf{P}}
\DeclareMathOperator{\E}{\textsf{E}\,}
\DeclareMathOperator{\var}{\textrm{var}}
\DeclareMathOperator{\sd}{\textrm{sd}}

%%% Příkaz pro transpozici vektoru/matice
\newcommand{\T}[1]{#1^\top}

%%% Vychytávky pro matematiku
\newcommand{\goto}{\rightarrow}
\newcommand{\gotop}{\stackrel{P}{\longrightarrow}}
\newcommand{\maon}[1]{o(n^{#1})}
\newcommand{\abs}[1]{\left|{#1}\right|}
\newcommand{\dint}{\int_0^\tau\!\!\int_0^\tau}
\newcommand{\isqr}[1]{\frac{1}{\sqrt{#1}}}

%%% Vychytávky pro tabulky
\newcommand{\pulrad}[1]{\raisebox{1.5ex}[0pt]{#1}}
\newcommand{\mc}[1]{\multicolumn{1}{c}{#1}}


\overfullrule=0pt

%% Titulní strana a různé povinné informační strany
\begin{document}
%%% Titulní strana práce a další povinné informační strany

%%% Titulní strana práce

\pagestyle{empty}
\hypersetup{pageanchor=false}

\begin{center}

\centerline{\mbox{\includegraphics[width=166mm]{../img/logo-cs.pdf}}}

\vspace{-8mm}
\vfill

{\bf\Large BAKALÁŘSKÁ PRÁCE}

\vfill

{\LARGE\AutorPrace}

\vspace{15mm}

{\LARGE\bfseries\NazevPrace}

\vfill

\Katedra

\vfill

{
\centerline{\vbox{\halign{\hbox to 0.45\hsize{\hfil #}&\hskip 0.5em\parbox[t]{0.45\hsize}{\raggedright #}\cr
Vedoucí bakalářské práce:&\Vedouci \cr
\noalign{\vspace{2mm}}
Studijní program:&\StudijniProgram \cr
\noalign{\vspace{2mm}}
Studijní obor:&\StudijniObor \cr
}}}}

\vfill

% Zde doplňte rok
Praha \RokOdevzdani

\end{center}

\newpage

%%% Následuje vevázaný list -- kopie podepsaného "Zadání bakalářské práce".
%%% Toto zadání NENÍ součástí elektronické verze práce, nescanovat.

%%% Strana s čestným prohlášením k bakalářské práci

\openright
\hypersetup{pageanchor=true}
\pagestyle{plain}
\pagenumbering{roman}
\vglue 0pt plus 1fill

\noindent
Prohlašuji, že jsem tuto bakalářskou práci vypracoval(a) samostatně a výhradně
s~použitím citovaných pramenů, literatury a dalších odborných zdrojů.
Tato práce nebyla využita k získání jiného nebo stejného titulu.

\medskip\noindent
Beru na~vědomí, že se na moji práci vztahují práva a povinnosti vyplývající
ze zákona č. 121/2000 Sb., autorského zákona v~platném znění, zejména skutečnost,
že Univerzita Karlova má právo na~uzavření licenční smlouvy o~užití této
práce jako školního díla podle §60 odst. 1 autorského zákona.

\vspace{10mm}

\hbox{\hbox to 0.5\hsize{%
V \hbox to 6em{\dotfill} dne \hbox to 6em{\dotfill}
\hss}\hbox to 0.5\hsize{\dotfill\quad}}
\smallskip
\hbox{\hbox to 0.5\hsize{}\hbox to 0.5\hsize{\hfil Podpis autora\hfil}}

\vspace{20mm}
\newpage

%%% Poděkování

\openright

\noindent
\Podekovani

\newpage

%%% Povinná informační strana bakalářské práce

\openright

\vbox to 0.5\vsize{
\setlength\parindent{0mm}
\setlength\parskip{5mm}

Název práce:
\NazevPrace

Autor:
\AutorPrace

\TypPracoviste:
\Katedra

Vedoucí bakalářské práce:
\Vedouci, \KatedraVedouciho

Abstrakt:
\Abstrakt

Klíčová slova:
\KlicovaSlova

\vss}\nobreak\vbox to 0.49\vsize{
\setlength\parindent{0mm}
\setlength\parskip{5mm}

Title:
\NazevPraceEN

Author:
\AutorPrace

\TypPracovisteEN:
\KatedraEN

Supervisor:
\Vedouci, \KatedraVedoucihoEN

Abstract:
\AbstraktEN

Keywords:
\KlicovaSlovaEN

\vss}

\newpage

\openright
\pagestyle{plain}
\pagenumbering{arabic}
\setcounter{page}{1}


%%% Strana s automaticky generovaným obsahem bakalářské práce

\tableofcontents

%%% Jednotlivé kapitoly práce jsou pro přehlednost uloženy v samostatných souborech
\chapter*{Úvod}
\addcontentsline{toc}{chapter}{Úvod}

Při trénincích a~soutěžích ve skoku o~tyči využívají atleti spolu s~trenéry videozáznamy provedených skoků. Videozáznam poté využívají pro detekci chyb v~technickém provedení skoku a~jejich případný vliv na výsledný výkon. Pro výkon při skoku o~tyči jsou kromě techniky důležité take biomechanické parametry skoku. Mezi tyto parametry patří napřáklad rychlost atleta při rozběhu.

V~tréninkovém procesu je důležité sledovat výkonnost atleta. K~tomuto účelu slouží především měření času běhu různých vzdáleností. Pro skok o~tyči je vhodné aplikovat podobné metody právě při skokanských trénincích. Jednou z~možností, jak toho docílit, je použití specializované techniky pro měření rychlosti atleta při rozběhu. Mezi specializovanou techniku lze zařadit radary nebo fotobuňky. Za použití vhodně rozmístěných fotobuněk lze poměrně přesně určit rychlost v~různých částech rozběhu. Využití těchto metod je finančně náročné a~jediným parametrem, který lze sledovat je pouze rozběhová rychlost atleta.

S~rozvojem počítačového vidění přibývá na světových atletických soutěžích studií zabývajících se pohyby sportovců \citep{IAAFBR}. Výsledky těchto studií jsou pro atlety a~jejich trenéry vynikající zpětnou vazbou. Tyto studie využívají několik kamer k~zachycení pohybu atleta, kamery nejprve kalibrují za použití přesně změřených konstrukcí s~vyznačenými body. Pro analýzu pohybu poté využívají 3D model atleta, který z~videí získají. Výsledný model zachycuje polohu kloubů těla a~na základě jejich pohybů lze jeho pohyb analyzovat. Takovýto model je možné použít k~analýze mnohem více parametrů skoku o~tyči než jen rychlosti běhu atleta.

Studie se většinou zabývají porovnáváním pohybu různých atletů na základě získaných parametrů. V~tréninkovém prostředí je zvykem využívat pouze videozáznamy a~případné porovnání skoků probíhá spuštěním záznamů, které dané skoky zachycují, následným rozborem techniky jednotlivých skoků a~hledáním případných rozdílů. Tato forma porovnání skoků ovšem nebere v~potaz biomechanické aspekty pohybu atleta, které mohou techniku skoku ovlivnit. Hodnoty parametrů, kterými se specializované studie zabývají, lze z~videa odhadnout pro účely porovnání skoků, ale ani zkušení trenéři nejsou schopni přesně určit hodnoty některých parametrů.

Podobné studie pohybu jsou užitečným tréninkovým prostředkem, ovšem jejich využití byť jen v~malém počtu tréninků se nejeví jako reálné. Dostupnějším přístupem k tomuto problému by byla aplikace, běžící na mobilním zařízení, která by analyzovala pohyb atleta na základě 3D modelu získaného z~jediného videa. Výsledkem této analýzy by byly graficky znázorněné parametry skoku, které by se daly využít pro porovnání různých skoků. Takováto aplikace by výrazně přispěla k~přesnější analýze tréninkových a~závodních skoků mimo soutěže světové úrovně a~pomohla by atletům lépe sledovat vlastní progres z~fyzické stránky.

Cílem mé práce je základní funkcionalita, kterou by takováto aplikace disponovala. Jedná se konkrétně o~extrakci modelu atleta z~jediného videozáznamu, detekce důležitých momentů skoku, mezi něž patří začátek rozběhu, odraz a~kulminace nad laťkou. Výsledný model poté analyzuji pro zisk hodnot biomechanických parametrů, které je při skoku vhodné zkoumat, tyto parametry poté zobrazím uživateli.

Text je rozdělen do čtyř kapitol. První kapitola se věnuje teorii skoku o~tyči a~počítačovému vidění. Ve druhé kapitole popisuji obecný návrh mého řešení. Třetí kapitola popisuje konkrétní implementaci a~ve čtvrté, závěrečné, kapitole se věnuji experimentům, které popisují funcionalitu programu.



\section*{Související práce}
\addcontentsline{toc}{section}{Související práce}



\chapter{Teorie}

\section{Skok o~tyči}

Skok o~tyči je jedna z~technicky náročných atletických disciplín. Atlet se nejprve s~tyčí rozběhne po rozběžisti, následně zasune tyč do kastlíku, poté se odrazí, provede skok a dopadne do doskočiště. Kastlík je místo v~zemi, kam atlet při rozběhu zasune tyč, aby měl při skoku stabilní oporu.

Atlet se při skoku pohybuje vzhůru za pomoci narovnávající se tyče, kterou ohne s~použitím energie získané při rozběhu. Tyče se liší délkou a~tvrdostí, tvrdší tyč je náročnější ohnout, aby atleta vynesla přes laťku, ale akumulovaná energie tyče je vyšší, a~tedy vyústí ve vyšší skok.

Pro pozici těla atleta při rozběhové fázi je typický značně omezený pohyb paží a~mírné natočení trupu. Důvodem tohoto nestandardního způsobu běhu je skutečnost, že atlet nese tyč, a~přesto se snaží vyvinout maximální kontrolovanou rychlost. Rychlosti, které atlet dosahuje v~momentu odrazu, se říká náběhová rychlost. Při skoku o~tyči (a~všech ostatních skokanských disciplínách) má atlet vyznačené místo, ve kterém začíná svůj pokus. Toto místo si volí sám, vliv na jeho polohu má především fyzická zdatnost atleta a~počet kroků rozběhu. Počet kroků závodního rozběhu atletů světové úrovně se pohybuje v~rozmezí 16 až 20~kroků, tedy přibližně 35 až 45~metrů od zadní hrany kastlíku. Atlet se rozbíhá s~tyčí ve vzduchu, postupně ji spouští, dokud není rovnoběžně se zemí. Následně provádí zásun - pohyb, při němž zasune tyč do kastlíku a~dostane paže nad hlavu. V~pozici s~pažemi nad hlavou se atlet odráží a přenáší energii do tyče. Na charakter skoku má vliv také místo odrazu. Následný skok lze rozdělit do několika částí.

První z~nich je odraz, při němž je kladen důraz na efektivitu přenosu energie získané při rozběhu do tyče. Po této krátké fázi skoku následuje zvrat. Jedná se o~pohyb způsobený švihnutím odrazové nohy a~paže, která se tyče drží výše, proti sobě - směrem dopředu. Po provedení zvratu se atlet dostane do pozice vzhůru nohama. V~této pozici atlet provádí obrat, při němž se otáčí o~180~stupňů podél vertikální osy, aby byl čelem k laťce. Následuje odraz od tyče a~přechod laťky, v~této fázi se atlet snaží dostat boky co nejvýš a~zajisit tak plynulý skok přes laťku. Laťku atlet překonává nohama napřed, čelem k~laťce.

\subsection{Analýza technického provedení skoku}

Způsob provedení skoku je podstatným ukazatelem pro výsledný výkon. Jen drobná změna v~jediné fázi pokusu může dramaticky ovlivnit charakter a~výšku celého skoku. Tím pádem mají skokané o~tyči pro nácvik techniky vyhrazeno několik tréninků týdně. Jedná se o~skokanské tréninky a~tréniny zabývající se nácvikem techniky prostřednictvím gymnastických prvků simulujících pohyb atleta na tyči. Pro mou práci jsou důležitější skokanské tréninky.

Skokanské tréninky začínají podobně, jako ostatní atletické tréninky, tedy rozcvičením. Po klasickém rozcvičení následuje příprava na samotné skákání, tato příprava se u jednotlivých atletů může lišit. Mnoho atletů před samotným skákáním provede několik cvičných zásunů a~rozběhů s~tyčí. Následují skoky z~krátkého rozběhu, standardně se jedná o 6 až 8~kroků, někteří atleti provádí tyto skoky bez ohýbání tyče. Již na základě těchto cviků lze určit, na jaké prvky skoku by se měl atlet zaměřit. Po této fázi následuje skákání z~dlouhého rozběhu. Počet kroků se liší podle fáze sezony, ve které se atlet nachází. V~závodním období jsou typické rozběhy delší, v~přípravném období kratší, jelikož tréninky v~přípravném období jsou zpravidla fyzicky náročnější.

Při tréninku se opakuje následující situace. Atlet provede skok a~následně ho konzultuje s~trenérem. Předmětem konzultace je především technické provedení skoku. S~rozvojem moderních technologií jsou konzultace ve většině případů doprovázeny sledováním a~rozborem videozáznamu zahycujícího právě provedený skok. S~pomocí tohoto videozáznamu lze přesně určit místo odrazu a~pohyb atleta při rozběhu a~následném skoku. Na pořízeném videozáznamu lze spolehlivě detekovat technické nedostatky skoků, které by se atlet měl v~následujících pokusech snažit eliminovat.

Videozáznam je nejčastěji pořizován na mobilní telefon, případně tablet. Pozice kamery se nejčastěji nachází na kolmici k~rozběžišti, která prochází místem odrazu. Z~tohoto místa je poměrně dobře vidět jak rozběh, tak skok. Navíc se takto s~velikou přeností dá určit místo odrazu a~pozice atletova těla při přenosu energie do tyče při odrazu. Strana, z~které je skok natočen se často mění, zálaží na prvku skoku, který má trenér s~atletem v~plánu zkoumat.

Závody probíhají podobně, po každém skoku opět dochází ke konzultaci, ovšem na mezinárodních závodech není běžné, že by atlet viděl záznam skoku, který trenér pořídil, na většině soutěží je to zakázané pravidly. Trenér tedy pouze popisuje nedostatky skoku a~probíhá diskuse s~atletem ohledně zaměřující se následujícím postupem. Při závodech není kladen takový důraz na změny v~technice jako při tréninku, typicky probíhá jen rozbor detailů nebo posouvání místa odrazu pro optimální přenos energie do tyče.

\subsection{Analýza biomechanických parametrů skoku}
\label{ssec:parametry}

S~rozvojem moderní techniky přibývá studií zabývajících se pohybem těla atleta při skoku o~tyči. Podkladem pro tyto studie jsou především mezinárodní závody, případně mistrovství republiky. Příkladem je biomechanická zpráva z~finále mužů Mistrovství světa v~atletice 2017 \citep{IAAF2017}, na níž se podíleli pracovníci Leeds Beckett University.

K~analýze parametrů se využívá několik kamer, jejichž záznam slouží jako podklad k~vytvoření modelu atleta v průběhu celého skoku. Tento model je následně analyzován pro zisk konkrétních parametrů skoku. Podobné studie se nejčastěji zabývají následujícími parametry:

\subsubsection{Délka jednotlivých kroků rozběhu}

Na efektivitu přenosu energie do tyče má vliv délka kroků na konci rozběhu. U~většiny atletů je poslední krok rozběhu kratší oproti předešlým krokům. Rozdíl v~délce předposledního a~posledního kroku se může pohybovat v~řádech desítek centimetrů. Důvodem pro zkrácení posledního kroku je sešlápnutí posledního kroku pod tělo atleta, tím atlet ztratí méně horizontální rychlosti při odrazu. Zkrácení posledního kroku doprovází odraz pod menším úhlem, což pro skok o~tyči není tak velký problém jako pro skokany do dálky.

\subsubsection{Doba oporové fáze kroku}

Technika běhu je značně ovlivněna dobou oporové fáze jednotlivých kroků. Doba oporové fáze souvisí s~rychlostí běhu, mírou pokrčení stojné nohy a výškou boků a kolen. Z~prostého videozáznamu se složitě určuje doba oporové fáze kroku, a~tudíž se při této analýze věnuje zvýšená pozornost právě míře pokrčení stojné nohy a výšce boků a kolen.

Doba oporové fáze kroku a~délka kroku se kromě analýzy videozáznamu dá měřit speciálními pásy, které se položí na zem po obou stranách rozběžiště.

\subsubsection{Náběhová rychlost}

Náběhová rychlost má značný vliv na přenos energie rozběhu do tyče, rychlejší atleti tedy zpravidla používají tvrdší a~delší tyče, které umožňují vyšší skoky.

Náběhovou rychlost lze měřit několika způsoby. Pro tréninkové účely se nejčastěji používají rovnoměrně rozmístěné fotobuňky, z~doby běhu mezi nimi lze snadno vypočítat průměrnou rychlost v~daném úseku. Tuto metodu používají především sprinteři, pro skok o~tyči není příliš vhodná, rychlost se v~momentu odrazu dramaticky mění a~tyč může ovlivnit detekci fotobuňek. Následující možností je použití radaru, s~kterým lze zanést aktuální rychlost do grafu. Při snímání atleta ze zadu téměř nedochází k~chybám měření. Způsob měření náběhové rychlosti použitý ve výše zmíněné studii \citep{IAAF2017} spočívá v~analýze detekovaného modelu v~souřadném systému, který umožňuje výpočet rychlosti jednotlivých částí těla ve standardních jednotkách, typicky se uvádí v ms\textsuperscript{-1}.

\subsubsection{Výška boků v~průběhu rozběhu}

V~průběhu rozběhu je zajímavé pozorovat výšku boků, výrazné výkyvy výšky boků mají za následek výkyvy rychlosti. Plynulost rozběhu tedy koreluje s~plynulostí výšky boků v~průběhu rozběhu. Pro rovný skok je důležité, aby se atlet rozeběhl kontrolovaně, což se může projevit ve výšce boků, výkyvy ve výšce boků také souvisí s~proměnlivou délkou kroků a~dobou oporové fáze kroku. Nejdůležitějsí moment pro zkoumání výšky boků jsou poslední kroky rozběhu.

\subsubsection{Místo odrazu}

Vzdálenost místa odrazu od zadní hrany kastlíku je podstatným ukazatelem efektivity přenosu energie do tyče. Je vhodné, aby měl atlet v~místě odrazu horní paži přímo nad sebou. Vhodná vzdálenost místa odrazu je individuální, někteří atleti preferují odraz blíže ke kastlíku než jiní. Pozice odrazu se odvíjí od technického provedení skoku, výšky a~odrazových schopností atleta.

\subsubsection{Výška úchopu}

Hodnota tohoto parametru souvisí s~délkou tyče a~odvíjí se od ní místo odrazu. Může se uvádět jako vzdálenost úchopu horní ruky od spodního konce tyče nebo jako výška úchopu horní ruky nad zemí, když je tyč svislá v~kastlíku. Kastlík je hluboký 20~cm, tedy hodnota po měření první metodou je o 20~cm vyšší než s použitím druhé metody. Hodnota tohoto parametru má vliv na dobu trvání skoku, vyšší úchup znamená delší dobu trvání skoku.

\subsubsection{Úhel odrazu}

Pro správný převod energie rozběhu do tyče je užitečné zkoumat úhel odrazu. Uvádí se jako úhel mezi zemí a~směrem, kterým se po odrazu pohybuje těžiště. Větší úhel znamená vyšší odraz, ale také větší ztrátu rychlosti při odrazu. Optimální hodnota úhlu odrazu různých atletů se liší například v~závoslosti na jejich tělesné výšce.

\subsubsection{Úhly v~kloubech při odrazu}

Pozice těla při odrazu je pro přenos energie do tyče také důležitá. Pokud je atletův trup nakloněn dopředu, jeho ramena a~paže jsou v~lepší pozici pro roztlačení tyče. Naopak pokud je atlet v~záklonu, jeho boky po odrazu snáze ujedou směrem dopředu a~ramena zůstanou vzadu. To má za důsledek horší roztlačení tyče.

\subsubsection{Doba trvání skoku}

Doba trvání skoku je vymezena odrazem od země a~opuštěním tyče. Na hodnotu tohoto parametru má vliv mnoho faktorů. Technika je jedním z~nich, standardně se atleti odráží od tyče vzpaženou paží, jsou ovšem atleti, kteří se tyče pustí dříve. Doba skoku za použití klasického odražení se vzpaženou paží se pohybuje okolo 1,3~s, nestandardní technika může tuto hodnotu stlačit až k 0,7~s. Tyto hodnoty se týkají mužského finále z~Mistrovství světa v~atletice 2017. Dále má na délku skoku vliv tvrdost a~délka tyče, tvrdší tyč se snáze - a~tedy i~rychleji - narovná a~skok tak trvá kratší dobu. S~použitím kratší tyče atletovo tělo překoná kratší vzdálenost a~doba trvání skoku je tedy kratší.

\subsubsection{Převýšení}

S~opuštěním tyče souvisí i~převýšení. Jedná se o~rozdíl zdolané výšky a~výšky úchopu nad zemí. Atleti světové úrovně převýšují okolo 130~cm, ženy okolo 70~cm. Převýšení je jedním z~hlavních parametrů, které určují výkonnost atleta. Na jeho hodnotu má vliv jak fyzická zdatnost, tak technické provedení skoku.




\section{Zpracování obrazu}



\subsection{Detekce člověka}

Postup detekce konkrétního objektu ve snímku se skládá ze dvou kroků. Prvním je zvolení vhodné reprezentace snímku (extrakce jeho vlastností), druhým krokem je spuštění klasifikátoru, který určí, zda (a~kde) se ve snímku hledaný objekt nachází.

Pro účely detekce se využívají okna, která reprezentují část snímku, na které se spustí samotná detekce. Pozice okna udá případnou pozici objektu ve snímku. Příklad tohoto využití je vidět na obrázku \ref{fig:bbox}, detekce postavy se spouští na různě posunutých oknech, které určují posunuté ohraničující rámečky.

Velikost okna se pro detekci různě velkých objektů mění, spíše než okno se mění velikost celého snímku, aby byla reprezentace různých částí snímku stejně velká. Touto modifikací je možné detekovat různě velké objekty v~tomtéž snímku.


\subsubsection{Extrakce vlastností snímku}

Velice úspěšným způsobem extrakce vlastnotí okna pro detekci postav je určení histogramu orientovaných gradientů \citep{HOG}. Tato metoda spočívá v~rozdělení okna na části (buňky), v~nichž se určí histogram, popisující směr gradientů - a~tedy i~směr hran. Jelikož gradient závisí na kontrastu snímku, části snímku je tedy vhodné normalizovat, co se týče kontrastu, aby byly jejich gradienty porovnatelné. Typicky se normalizují skupiny buňek (bloky), ne jednotlivé buňky oken. Získané hodnoty z~takto normalizovných bloků, které pokrývají okno a z~části se překrývají, reprezentují vlastnosti daného okna, které se následně předají klasifikátoru.

Výsledné vlastnosti jsou invariantní k~lokálním geometrickým transformacím, mezi které patří posuny, škálování a~rotace. Tyto transformace mají na výslednou reprezentaci objektu vliv v~případě, že jsou příliš velké. Díky této vlastnosti je detekce člověka úspěšná, pokud se nachází ve~víceméně vzpřímené pozici. O~předpoklad vzpřípené pozice těla se opírá většina detektorů.

Extrakce těchto vlastností dosahuje vysoké úspěšnosti při použití přístupu Scale-Invariant Feature Transform \citep{SIFT}. Tento přístup představuje invarianci vlastností vzhledem ke škálování a~rotaci.


\subsubsection{Klasifikace}

Při využití histogramu orientovaných gradientů se nejčastěji využívá lineární Support vector machines (SVM) \citep{SVM}. SVM se používá k~binární klasifikaci.

Základ metody SVM spočívá v~rozdělení prostoru na dva poloprostory s~využitím nadroviny. Při trénování mám daná data rozdělená do dvou tříd. Lineární SVM předpokládá lineárně separabilní data, tedy je lze oddělit nadrovinou. Nelineární SVM využívá k~dosažení stejného cíle prostorové transformace. Při trénování hledám nadrovinu, která data odděluje a~vzdálenost nadoviny od nejbližšího bodu z~trénovacích dat je maximální možná.

Každý bod v~prostoru odpovídá příslušné kombinaci vlastností vstupu. Podle vlastností vstupu určím pozici v~prostoru a~podle poloprostoru, do kterého spadá, vstup klasifikuji.

V~mém případě určím vlastnosti vstupního snímku, na základě těchto vlastností určím poloprostor, do něhož snímek spadá a~následně se rozhodnu, zda je na obrázku postava či nikoliv.



\subsection{Trasování objektů}

Trasování objektů ve videu je metoda, díky níž lze na základě pozice objektu ve snímku určit jeho pozici v~následujícím snímku. Tímto postupem lze objekt sledovat v~průběhu celého videa. Pro lepší přesnost detekce je vhodné stav objektu postupně aktualizovat, jelikož se v~průběhu videa může menit. Algoritmům, které trasování objektů zajišťují, se říka trackery.

Přístupů k~řešení tohoto problému je několik, lze je rozdělit do následujících kategorií:
\begin{itemize}
    \item feature based,
    \item segmentation based,
    \item estimation based a
    \item learning based.
\end{itemize}

Trasovací metody lze rozdělit také podle jiných kritérií, příkladem je rozdělení pouze do $2$\,\rm kategorií - využívající korelační filtr a~ostatní \citep{}.

\subsubsection{Feature based}

Tento přístup spočívá v~extrakci vlastností objektu. Mezi tyto vlastnosti se řadí například barva nebo textura. Na základě extrahovaných vlastností se v~následujícím snímku najde objekt, který má tyto vlastnosti nejpodobnější a~splňuje další podmínky, mezi které lze zařadit vzdálenost nalezeného objektu od jeho pozice v~předchozím snímku.

\subsubsection{Segmentation based}

Trasování založené na principu segmentace odděluje objekty od pozadí. Typicky je trasován objekt v~popředí videa, tudíž oddělení všech objeků v~popředí od pozadí zjednoduší problém trasování.

\subsubsection{Estimation based}

Metody založené na odhadech se opírají o~reprezentaci objektu stavovým vektorem. Tento vektor popisuje směr pohybu objektu, který se neustále upravuje. Trasování tímto způsobem lze rozdělit do dvou kroků - predikce a~aktualizace.

Predikce využívá stavový vektor k~určení pozice objektu v~následujícím snímku a~aktualizační část predikovanou pozici upraví na základě analýzy této části snímku podobným způsobem, jako při použití feature based trackeru.

\subsubsection{Learning based}

Tento přístup zahrnuje metody využití strojového učení. K~detekci obejktů v~následujícím snímku je potřeba extrahovat ze snímku vlastnosti a~na jejich základě trénovat tracker.

Využití konvolučních sítí dosahuje výborných výsledků, ve výsledcích soutěže The Visual Object Tracking VOT2017 challenge \citep{VOT2017} byly mezi nejpřesnějšími algoritmy využívající porovnávání tvarů za použití konvolučních sítí (CNN matching). Nejpřesnější trackery využívají korelační filtry.



\subsection{Odhad pozice těla}

Odhad pozice těla spočívá v~detekci kloubů těla, k~této detekci se využívá především hluboké učení. Hluboké učení využívá rozsáhlé neuronové sítě, které se skládají z~několika vrstev. První vrstva bývá označovaná jako vstupní, poslední jako výstupní. Vrstvy, které se nachází mezi nimi jsou označovány jako skryté. Každou vrstvu tvoří neurony, které při zpracování vstupu nabývají číselných hodnot. Přechody mezi vrstvami definují váhy, kterými se hodnoty příslušných neuronů násobí. Většina vrstev využívá aktivační funkci, která hodnoty neuronů upravuje než je předá další vrstvě. Aktivační funkce jsou typicky nelineární, mezi nejpoužívanější se řadí ReLU a~tanh, která vznikla modifikací funkce sigmoid, jež se využívala v~minulosti a~slouží k~binární klasifikaci na výstupní vrstvě. Výstupem sítě jsou hodnoty neuronů výstupní vrstvy, které nejčastěji reprezentují pravděpodobnost, že vstup spadá do třídy, kterou daný neuron reprezentuje. Aby se jednalo o~pravděpodobnostní distribuci, musí se hodnoty neuronů normovat. K~tomu slouží funkce softmax, což je rozšíření funkce sigmoid pro klasifikaci do několika tříd. Aktivační funkce popisuje následující tabulka.

\begin{alignat*}{2}
\label{eq:aktivace}
    &\text{ReLU:} \qquad &\text{ReLU}(x) &= max(0,x) \\
    &\text{sigmoid:} \qquad &\sigma(x) &= \frac{1}{1+e^{-x}} \\
    &\text{tanh:} \qquad &\text{tanh}(x) &= 2\sigma(2x) - 1 \\
    &\text{softmax:} \qquad &\text{softmax}(\mathbf{x}) &= \frac{\mathbf{x}}{\sum_i{e^{\mathbf{x}_i}}} \\
\end{alignat*}

Existují různé typy neuronových sítí podle typu vrstev, které obsahují. První sítě, které se využívaly, obsahovaly vrstvy, které byly kompletně propojené. Hodnoty všech neuronů dané vrstvy v~takové síti závisí na hodnotách všech neuronů předchozí vrstvy. Tyto sítě není příliš vhodné používat pro detekci objektů ve snímcích, jelikož závisí na mnoha vlastnostech množiny dat, na kterých jsou natrénovány. Hlavní nevýhoda spočívá v~tom, že tato síť není invariantní k~posunu objektu. Pokud je síť natrénovaná na obrázcích krajiny, kde se slunce nachází v~levém horním rohu, má poté problém detekovat slunce v~pravém horním rohu. Pixely, které reprezentují slunce se totiž nachází na jiném místě ve vstupní vrstvě a~jejich hodnoty využívají jiné váhy.

Řešení tohoto problému nabízí konvoluční sítě. Ty využívají konvoluční filtry, které jsou specifické pro každou vrstvu. Takovýto filtr má danou velikost, ze které bere vstupní hodnoty a~vypočítá hodnotu výslednou. Jedná se o~jakési okno, které se aplikuje na každou část vstupní vrstvy a určí hodnotu konvoluce v~této části. Hodnota neuronu vrstvy tedy závsí jen na blízkém okolí v~předchozí vrstvě. Jelikož se hodnoty filtru při zpracování dané vrstvy nemění, je výsledná hodnota po aplikaci filtru na části obsahující slunce v~levém horním rohu snímku stejná, jako výsledná hodnota po aplikaci na část obsahující slunce v~pravém horním rohu snímku. Více konvolučních vrstev dokáže popsat větší část snímku, jejíž velikost závisí na velikostech použitých filtrů.







































\chapter{Návrh řešení}

V~této kapitole stručně popíšu návrh mého programu, který má za úkol extrakci modelu reprezentujícího pohyb atleta při skoku o~tyči. Skok je programu předán ve formě videa, po extrakci modelu program detekuje důležité momenty skoku a~analyzuje vhodné biomechanické parametry. Hodnoty těchto parametrů následně zobrazí uživateli v~grafické podobě a~také zobrazí snímky zachycující důležité momenty skoku.

Program lze spustit ve dvou módech - analýza videa a~analýza modelu.

Při analýze videa dostane program na vstupu video, ve videu najde atleta a s~využitím jeho pozice detekuje jeho kostru v~jednotlivých snímcích. Z~detekované kostry vytvoří 3D model atletova těla, který reprezentuje atletův pohyb.

Analýza pohybu modelu probíhá stejně při zpracování videa i~samotného modelu. Podle pohybu jednotlivých částí těla atleta se vypočítají hodnoty biomechanických parametrů, které se uloží a~zobrazí uživateli.




\section{Souřadné systémy}
\label{sec:ssystemy}

Pro reprezentaci pozic objektů je potřeba zvolit vhodné souřadné systémy, aby nezkreslovaly hodnoty parametrů.



\subsection{Snímek}

Souřadný systém ve snímku má počátek v~levém horním rohu. Bod $(x,y)$ je ve vzdálenosti $x$\,\rm pixelů od levé strany snímku a $y$\,\rm pixelů od horní strany snímku.



\subsection{Video}

Kamera typicky není při natáčení skoků statická, tudíž je potřeba reprezentovat pohyb objektů jinak než pouze jejich pozicí ve snímcích videa.

Rozhodl jsem se pro určení pohybu podle vzájemné pozice objektu a~pozadí videa v~souřadném systému snímku. Polohu pozadí považuji za statickou. Jedná se o~jistou formu projekce 3D prostoru do 2D, která zkresluje realitu, tento systém lze reprezentovat jako souřadnice v~panoramatické fotce.

Pozice objektu v~prvním snímku, v~němž ho detekuji, je počátkem souřadného systému. Nejprve určím pozici levého horního rohu snímku, v~němž se nachází objekt, jehož pozici chci určit. Tuto pozici určím vzhledem k~pozici pozadí. Společně s~pozicí objektu v~tomto snímku získám pozici objektu vůči jeho pozici ve snímku, v~němž jsem ho detekoval poprvé.

Souřadný systém videa znázorňuje obrázek \ref{fig:ssystem_video}.

\begin{figure}[h]\centering
    \includegraphics[width=\textwidth]{ssystem_video}
    \caption{
        \centering\small
        Vizualizace souřadného systému videa. Světlé pole reprezentuje snímek videa, $p$ posun jeho horního rohu vůči původní pozici atleta, $q$ bod, jehož pozici ve videu získám pomocí posunu snímku a~pozice bodu $q$ ve snímku.
    }
    \label{fig:ssystemvideo}
\end{figure}




\subsection{Prostor}

Model atleta umístím do 3D prostoru, který se podobá reálnému světu, ale je založen na souřadném systému videa, tudíž realitu do jisté míry zkresluje.

Odhad reálných vzdáleností z~videa je poměrně komplikovaný a~pro základní parametry skoku bude tento systém dostačující. Pro případnou konverzi na metry lze systém vhodně přeškálovat v~budoucnu, zatím vzdálenosti odpovídají vzdálenostem ve snímku videa, které závisí na jeho rozlišení. Musí se ovšem brát v~potaz vzdálenost atleta od kamery, jelikož se jeho velikost v~průběhu videa mění. Pro přesnější konverzi bych tedy musel zvolit jiný způsob určení pohybu objektů ve videu.

Kloub výsledného 3D modelu reprezentuji jako uspořádanou trojici $(x,y,z)$, počátek této soustavy umístím do místa kotníku odrazové nohy při odrazu atleta. Ideální by bylo počátek umístit na zadní hranu zasouvací skříňky. Zasovací skříňka je schovaná za doskočištěm, tudíž bych jeho pozici musel odhadovat na základě detekce tyče, což by mohlo být nepřesné.

Hodnota první složky určuje horizontální vzdálenost bodu od počátku ve směru rovnoběžném s~osou rozběžiště. Tato hodnota stoupá ve směru rozběhu.

Hodnota druhé složky určuje horizontální vzdálenost bodu od počátku ve směru kolmém na osu rozběžiště. Tato hodnota stoupá ve směru doprava z~pohledu atleta při rozběhu.

Hodnota třetí složky určuje vertikální vzdálenost bodu od počátku. Tato hodnota stoupá ve směru vzhůru.

Program zatím využívá dvoudimenzionální podprostor 3D prostoru z~důvodu komplikovaného odhadu hodnoty druhé složky pozice zkoumaných bodů. 3D prostor jsem implementoval s~výhledem budoucího rozšíření aplikace.

Souřadný systém videa je znázorněný na obrázku \ref{fig:ssystemprostor}.

\begin{figure}[h]\centering
    \includegraphics[width=\textwidth]{ssystem_prostor}
    \caption{
        \centering\small
        Vizualizace souřadného systému 3D prostoru. Červená reprezentuje osu $x$, zelená $y$ a~modrá $z$. Směr šipek odpovídá růstu hodnot bodů v~prostoru.
    }
    \label{fig:ssystemprostor}
\end{figure}




\section{Vstup programu}
\label{sec:vstup}

Vstupem programu je video skoku o~tyči pořízené běžnými prostředky, jakými jsou mobilní telefon nebo tablet.

Videa skoku o~tyči se nejčastěji natáčí ze strany. Kamera se standardně nachází na kolmici k~rozběžišti, která prochází místem odrazu. Můj program očekává na vstupu video ze strany, pro přesnost hodnot výsledných parametrů je vhodné, aby byla kamera na úrovni odrazu. Pro potřebu přesnější analýzy běhu atleta je lepší umístit kameru na úroveň poloviny rozběhu.

Při trénincích se kamera nachází ve vzdálenosti okolo 10 metrů od rozběžiště, závodní skoky jsou ve většině případů natáčeny z~tribuny, tudíž je vzdálenost větší (v~řádu desítek metrů). Velikost atleta ve videu nemá na funkcionalitu programu vliv, ale detekce menší postavy může vyústit v~horší přesnost.




\section{Analýza videa}

Vstupem programu je video skoku o~tyči. Toto video program nejprve zpracuje do vhodné reprezentace, následně ve videu nalezne atleta. Na základě pozice atleta v~jednotlivých snímcích detekuje kostru atleta, kterou převede do 3D modelu. Podle pohybu atletova těla, který je reprezentovaný získaným modelem, analyzuje parametry skoku a~určí jejich hodnoty. Výstupem programu je model atletova těla a~hodnoty definovaných parametrů. Tento výstup následně uloží a~zobrazí uživateli.

Program lze spustit na počítači s~Unixovým operačním systémem.



\subsection{Zpracování videa}

Prvním krokem k~provedení analýzy je zpracování videa. Aby video nebylo otevřené po celou dobu analýzy, rozhodl jsem se při analýze pracovat s~jednotlivými snímky, které si ukládám při jediném průchodu videem. Rozlišení videa může být libovolné, ale z~důvodu časové náročnosti rozlišení snímků videa před analýzou zmenším.



\subsection{Nalezení atleta}

Pro nalezení atleta stačí určit snímek, ve kterém je atlet dobře vidět, a~jeho pozici v~něm. Pozici atleta bude reprezentovat jeho ohraničující rámeček, který je vyznačený na obrázku \ref{fig:bbox}.

\begin{figure}[h]\centering
    \includegraphics[width=\textwidth]{bbox}
    \caption{
        \centering\small
        Ohraničující rámeček atletova těla ve snímku videa \texttt{1.MOV} (viz tabulka \ref{fig:videa}).
    }
    \label{fig:bbox}
\end{figure}

Pro nalezení atleta jsem zvolil dva přístupy.

Prvním z~nich je manuální. Uživatel po spuštění analýzy vybere snímek, ve kterém je atlet dobře vidět. Následně vyznačí do videa ohraničující rámeček, do něhož se vejde postava atleta.

Druhým přístupem je automatické nalezení atleta. Při automatické detekci atleta je potřeba detekovat postavy ve snímku videa. Výsledkem této detekce je množina ohraničujících rámečků reprezentujících postavy ve snímku. Z~těchto postav je potřeba vybrat tu, která reprezentuje atleta.

Atlet provádí na videu specifický pohyb. Nejprve se rozeběhne a~poté provede skok, tedy ohraničující rámeček postavy atleta by se ve videu měl pohybovat horizontálně bez výrazných výkyvů ve vertikálním směru, následně se pohybovat směrem vzhůru a~nakonec směrem dolů.

Pomocí trasování objektů mohu ve videu zkoumat pohyb postav a~na základě jejich reálného pohybu filtrovat postavy, které mohou reprezentovat atleta.

V~prvním snímku detekuji postavy a~uložím si je do seznamu, následně budu zkoumat jejich pohyb. Jejich pohyb budu zkoumat pomocí trasování objektů v~souřadném systému videa.

Pokud se postava v~horizontálním směru nepohne za určitou dobu alespoň o~danou vzdálenost, smažu ji ze seznamu, jelikož nereprezentuje atleta. Jakmile je seznam postav prázdný, spustím detekci postav znovu v~právě zpracovávaném snímku.

V opačném případě považuji postavu za atleta. Tímto okamžikem hledání atleta ve videu končí, jelikož vím, ve kterém snímku jsem atleta detekoval poprvé a~jaký byl jeho ohraničující rámeček v~daném snímku.


\subsubsection{Trasování postav}

Trasování postav ve videu používám k~odlišení atleta od ostatních postav na základě analýzy jejich pohybu. Trasování atleta provádím také při detekci kloubů kostry, ale modifikovaným způsobem, který je popsaný v~sekci \ref{ssec:ndetekce}.



\subsection{Detekce kostry}
\label{ssec:ndetekce}

\subsubsection{Trasování atleta}

Trasování atleta při rozběhu nedělá většině trackerů problém. Nepřesné trasování nastává při skoku. Důvodem těchto nepřesností je rotace atletova těla. Tato rotace je při skoku výrazná a~při vodorovné pozici trupu nemusí atletovo tělo zabírat věšinu ohraničujícího rámečku, což způsobí, že tracker začne trasovat pozadí videa a~ztratí atleta. Tato vlastnost a~její napravení je vidět na obrázku \ref{fig:ztrata}.

\begin{figure}[p]\centering
    \subfloat{
        \includegraphics[width=.45\textwidth]{rozbeh_bbox}
    }
    \subfloat{
        \includegraphics[width=.45\textwidth]{rozbeh_grid}
    } \\
    \subfloat{
        \includegraphics[width=.45\textwidth]{odraz_bbox}
    }
    \subfloat{
        \includegraphics[width=.45\textwidth]{odraz_grid}
    } \\
    \subfloat{
        \includegraphics[width=.45\textwidth]{zvrat_bbox}
    }
    \subfloat{
        \includegraphics[width=.45\textwidth]{zvrat_grid}
    } \\
    \subfloat{
        \includegraphics[width=.45\textwidth]{skok_bbox}
    }
    \subfloat{
        \includegraphics[width=.45\textwidth]{skok_grid}
    }
    \caption{
        \centering\small
        Použití jednoho trackeru (vlevo) a~mřížky trackerů, která je implementovaná v~programu (vpravo). Použité snímky jsou vybrané z~videa \texttt{33.MOV} (viz tabulka \ref{fig:videa}).
    }
    \label{fig:ztrata}
\end{figure}

Modifikací vedoucích ke zlepšení přesnosti jsem zkusil několik.

Nejprve jsem otáčel snímky videa podle naklonění trupu v~předešlém snímku. K~určení tohoto naklonění jsem použil detekovanou kostru atletova těla.

Následně jsem aktualizoval tracker, aby trasoval pouze trup. Jeho pozici jsem opět získával z~detekované kostry. Tyto metody selhávaly při chybné detekci kostry, ačkoliv jsem se snažil implementovat kontrolní mechanismus. Například jsem neaktualizoval tracker na kostře, jejíž trup byl příliš daleko od právě trasovaného trupu. Problém detekce kostry špatné postavy je vidět na obrázku \ref{fig:spatnakostra}.

Nakonec jsem implementoval metodu, která nezávisí na detekci kostry atleta.

Atletovo tělo netrasuji jako celek, ale po částech. Původní ohraničující rámeček rozdělím na mřížku menších rámečků. Každý rámeček trasuji zvlášť a~průběžně aktualizuji ty, které atleta ztratí.

Po každých několika snímcích kontroluji pohyb jednotlivých rámečků ve videu od předešlé kontroly. Určím rámeček, který se pohnul nejvíce. Pohyb rámečku musí být horizontálně ve stejném směru, jako je rozběh atleta. Následně posunu původní mřížku tak, aby odpovídala pozici nejvíce se pohybujícího rámečku. Do mřížky poté přesunu rámečky, jejichž trasování selhalo, na své původní místo. Průběh aktualizace je zobrazen na obrázku \ref{fig:mrizka}.

\begin{figure}[p]\centering
    \subfloat[]{
        \includegraphics[width=.45\textwidth]{mrizka1}
    }
    \subfloat[]{
        \includegraphics[width=.45\textwidth]{mrizka2}
    } \\
    \subfloat[]{
        \includegraphics[width=.45\textwidth]{mrizka3}
    }
    \subfloat[]{
        \includegraphics[width=.45\textwidth]{mrizka4}
    }
    \caption{
        \centering\small
        Pohyb mřížky v~průběhu její první aktualizace. (a)~-~úvodní pozice mřížky, (b)~-~aktualizace stále neproběhla, protože není znám směr rozběhu, (c)~-~poslední snímek před aktualizací mřížky, (d)~-~první snímek po aktualizaci mřížky
    }
    \label{fig:mrizka}
\end{figure}



\subsubsection{Detekce kloubů těla}

Pro určení atletovy kostry stačí detekovat klouby, které následně spojím úsečkami. Klouby detekuji pomocí konvoluční sítě. Jelikož má implementace očekává, že na snímku je jediná postava, bylo pro dostatečnou kvalitu detekce kloubů potřeba vyříznout ze snímku okno, v~němž se nachází atlet a~příslušně okno rotovat, aby byl atlet v~co nejvzpřímenější pozici. Síť, kterou používám, je totiž natrénovaná na databázi MPII Human Pose dataset \citep{MPIIHPE}, v~níž je většina postav ve vzpřímené poloze.

Určení pozice atletova okna provádím s~použitím rámečků na těle atleta, kterého trasuji. Rámečky uzavřu do co nejmenšího rámečku a v~něm naleznu střed. To je střed prvního okna pro detekci kostry. Střed okna v~následujících snímcích posouvám směrem ke kyčlím atleta. Tento posun získám z~pozice rámečku obsahujícího trackery atletova těla a~pozice kyčlí kostry detekované v~předchozím snímku. Aby se okno neposunulo příliš daleko od těla atleta při chybné detekci kostry, posouvám střed okna jen uvnitř rámečku obsahujícího trackery atletova těla. Tím napravím detekci kostry chybné postavy po několika snímcích, viz obrázek \ref{fig:spatnakostra}

\begin{figure}[p]\centering
    \subfloat[]{
        \includegraphics[width=.45\textwidth]{spatnakostra1}
    }
    \subfloat[]{
        \includegraphics[width=.45\textwidth]{spatnakostra2}
    } \\
    \subfloat[]{
        \includegraphics[width=.45\textwidth]{spatnakostra3}
    }
    \subfloat[]{
        \includegraphics[width=.45\textwidth]{spatnakostra4}
    }
    \caption{
        \centering\small
        Detekce špatné kostry a~následné napravení. (a)~-~správná detekce, (b)~-~špatná detekce, (c)~-~průběh nápravy, (d)~-~správná detekce
    }
    \label{fig:spatnakostra}
\end{figure}

Rotaci okna provádím podle náklonu trupu poslední validně detekované kostry. Validně detekovanou kostrou rozumím kostru, které jsem detekoval všechny klouby. Abych zamezil rotacím podle chybné detekce, aktualizuji poslední známý úhel náklonu trupu jen v~případě, že se od předchozího příliš neliší. Pokud tedy při skoku detekuji kostru postavy, která stojí u~místa odrazu, neaktualizuji úhel rotace a~okno rotuji podle posledního známého naklonění trupu atleta. [OBRÁZEK]

Pokud se mi nepodaří detekovat celou kostru, zkusím okno rotovat o~určitý úhel oběma směry a~vyberu ten s~nejúspěšnější detekcí, co se počtů detekovaných kloubů kostry týče.

Jelikož atlet mění v~průběhu videa velikost - standardně se zvětšuje, protože se přibližuje ke kameře - je potřeba okno postupně zvětšovat. První okno je dané původním ohraničujícím rámečkem, který zvětším konstantou. Velikosti následujících oken určuji podle dosud největší validně detekované kostry. Velikost se rovná vzdálenosti dvou nejvzdálenějších kloubů kostry. Okno pro detekci kostry je čtvercové, jehož strana má délku velikosti dosud největší validně detekované kostry, jejíž velikost zvětším konstantou. Okno pro první detekci nemusí být čtvercové, odvíjí se od výsledku nalezení atleta ve videu.



\subsection{Konverze kostry do 3D modelu}

Konverzi detekované kostry do 3D modelu provedu pomocí algoritmu pro převod mezi souřadnými systémy videa a~3D modelu. Tyto souřadné systémy jsou popsané v~sekci \ref{sec:ssystemy} a~implementace algoritmu v~sekci \ref{sec:konverze}.



\subsection{Analýza parametrů}

Vzhledem k~vlastnostem mnou zvoleného modelu nebude vhodné analyzovat některé parametry skoku. Rozeberu proveditelnost analýzy biomechanických parametrů popsaných v~sekci \ref{ssec:parametry}.


\subsubsection{Délka jednotlivých kroků rozběhu}

Délka jednotlivých kroků rozběhu je velice citlivá na úhel, pod kterým atleta snímám, a~jeho vzdálenost od kamery. Vliv vzdálenosti by bylo možné vyřešit reprezentací délky kroku relativně vůči výšce postavy. Ovšem úhel osy kamery vůči rozběžišti v~programu nijak neřeším, tudíž se této nepřesnosti nelze jednoduše zbavit.

Proto jsem se rozhodl tento parametr neanalyzovat.

Místo délky kroku jsem se rozhodl implementovat zjišťování hodnoty doby trvání jednotlivých kroků. K~zisku této hodnoty stačí najít snímky, ve kterých se nižší noha začíná pohybovat směrem vzhůru. Tím získám (po jednoduché filtraci) momenty dokončení odrazů jednotlivých kroků. Doba, která uplyne mezi následujícími kroky, je výsledná hodnota.


\subsubsection{Doba oporové fáze kroků}

Dobu oporové fáze kroku lze analyzovat mírným rozšířením analýzy doby trvání jednotlivých kroků. Při spuštění stejného algoritmu od konce videa získám momenty došlapů. Spolu s~momenty odrazů lze určit dobu oporové fáze jednotlivých kroků.

Vzhledem k~tomu, že jeho přesnost závisí na frekvenci snímků videa, jsem se analýzu tohoto parametru rozhodl neanalyzovat.


\subsubsection{Náběhová rychlost}

Náběhovou rychlost nelze analyzovat v~průběhu celého rozběhu příliš přesně z~důvodu měnícího se úhlu, pod kterým kamera snímá atletův rozběh. V~momentu odrazu lze rychlost reprezentovat poměrně přesně, ale jen v~jednotkách závislých na pixelech. Pro převod na jednotky s~větší vypovídající hodnotou by bylo možné využít délku tyče nebo výšku atleta. Pro implementaci této metody by bylo potřeba, aby uživatel zadal příslušný parametr - výšku atleta nebo délku tyče - a~výsledná přesnost by nebyla příliš dobrá, jelikož se náběhová rychlost u~různých skoků liší nepatrně. Lepším způsobem pro analýzu tohoto parametru tak zůstává využití fotobuněk nebo radaru.

Místo analýzy tohoto parametru jsem zkoumal ztrátu horizontální rychlosti ramen a~boků při odrazu. Tato změna nezávisí na použitých jednotkách, takže není potřeba odhadovat vzdálenosti. Pro zisk těchto hodnot porovnávám změnu pozice ramen a~boků daný časový úsek před odrazem, při odrazu a~stejný časový úsek po něm.


\subsubsection{Výška boků v~průběhu rozběhu}

Výška boků bez jakéhokoliv škálování v~závislosti na velikosti atleta ve videu není příliš vypovídající pro porovnání začátku a~konce rozběhu. Vliv na výsledný skok mají spíše lokální výkyvy, především ty, které nastávají v~konci rozběhu. Tím pádem je hodnota tohoto parametru užitečná i~bez škálování naměřených hodnot.

Výšku boků jsem tedy mezi implementované parametry zahrnul, pro zisk jeho hodnoty průměruji výšku levé a~pravé kyčle.


\subsubsection{Místo odrazu}

Přesnost tohoto parametru závisí také na zvoleném modelu, kterým reprezentuji tělo atleta, tedy kolik bodů na těle detekuji. Nejčastěji se udává jako vzdálenost špičky chodidla od zadní hrany zasouvací skříňky. Pozici špičky chodidla je složité odhadnout, pokud znám jen pozici kotníku. Přesnost odhadu hrany zasouvací skříňky by na hodnotu tohoto parametru měla také značný vliv.

Vzdálenost místa odrazu se zkoumá kvůli jejímu vlivu na horizontální vzdálenost úchopu horní ruky od odrazové nohy. Tuto vzdálenost jsem schopen změřit poměrně přesně, bohužel pouze v~nedefinovaných jednotkách, případný odhad konkrétní vzdálenosti by nedosáhl větší přesnosti než odhad pouhým okem.

Pozice odrazu se jednoduše a~přesně určí pouhým okem, především pokud jsou poblíž místa odrazu na zemi značky, proto jsem tento parametr neimplementoval.


\subsubsection{Výška úchopu}

Podobně jako v~předchozím případě lze hodnotu tohoto parametru poměrně přesně určit pouhým okem při znalosti délky tyče. Proto nebylo potřeba analýzu tohoto parametru zatím implementovat.

V~budoucnu bych aplikaci rád rozšířil o~detekci tyče, která by možnosti analýzy skoku značně rozšířila.


\subsubsection{Úhel odrazu}

Jelikož je osa kamery při odrazu většinou kolmá na směr rozběhu, je možnost získání přesné hodnoty tohoto parametru solidní. Úhel určím podobně jako ztrátu rychlosti. Porovnám pozici boků daný časový úsek před odrazem, při odrazu a~stejný časový úsek po něm. Na základě rozdílů těchto pozic určím úhel rozběhu (jeho konce) a~skoku (jeho počátku) vůči horizontální ose. Po odečtení úhlu rozběhu od úhlu skoku získám úhel odrazu.

Úhel odrazu jsem se tedy rozhodl analyzovat.


\subsubsection{Úhly v~kloubech při odrazu}

Úhly v~kloubech kostry lze z~modelu získat snadno, ale lze je odhadnout poměrně přesně i~bez analýzy programem - alespoň pro účely porovnání skoků. Nejedná se o~tak důležitý parametr jako úhel odrazu, tedy implementaci analýzy úhlů kloubů těla ponechám jako možné rozšíření programu do budoucna.


\subsubsection{Doba trvání skoku}

Pro zisk hodnoty tohoto parametru by bylo užitečné detekovat tyč. Tuto funkcionalitu jsem neimplementoval, tudíž tento parametr není mezi analyzovanými.


\subsubsection{Převýšení}

Pro určení míry převýšení je také vhodné detekovat tyč. Druhou možností je převod souřadného systému na metry a~zadání výšky úchopu. Poté by bylo možné převýšení určit. V~budoucích verzích programu by se tento parametr mohl objevit. Zatím jsem ho neimplementoval.


\subsubsection{Další implementované parametry}

Nad rámec popsaných parametrů jsem se rozhodl implementovat analýzu následujících parametrů:

\paragraph{Úhel došlapu jednotlivých kroků}

Parametrem souvisejícím s~dobou oporové fáze kroku je úhel došlapu. Jedná se o~úhel mezi kotníkem, kyčlí a~vertikálou. Aproximuje vzdálenost došlapu před těžiště atleta. Pro zisk hodnoty tohoto parametru použiji metodu získávání snímků, ve kterých dochází k~došlapu.

\paragraph{Náklon trupu}

Jedná se o~náklon trupu v~rovině určené vertikálou a~osou rozběžiště. Hodnota tohoto parametru není příliš přesná na začátku rozběhu, ale s~postupem času se přesnost zvětšuje díky lepšímu úhlu záběru těla atleta.


\subsubsection{Důležité momenty skoku}

Kromě parametrů ve videu detekuji podstatné momenty skoku, za něž považuji začátek rozběhu, odraz a~moment kulminace boků nad laťkou.




\section{Analýza modelu}

Aby uživatel nemusel video analyzovat znovu, je mu umožněno načtení modelu ze souboru, který je popsaný v~sekci \ref{sec:vystup}.

Funkcionalita programu je velice podobná, na rozdíl od analýzy videa se na začátku načte také soubor s~uloženým modelem. Neprovádí se hledání atleta ve videu, detekce kostry, ani konverze kostry do 3D modelu. Analýza parametrů probíhá stejně jako při analýze videa, model se neukládá, ukládají se jen parametry, aby se jejich hodnoty daly využít pro případnou analýzu.

Kostru z~3D modelu získá program s~využitím konverzí souřadných systémů, kterou popisuji v~sekci \ref{sec:konverze}.




\section{Výstup programu}
\label{sec:vystup}

Výstupem programu je model reprezentující pohyb atleta v~průběhu rozběhu i~skoku a~hodnoty definovaných parametrů, které jsou
\begin{itemize}
    \item doba trvání jednotlivých kroků,
    \item ztráta rychlosti ramen a~boků,
    \item výška boků,
    \item úhel odrazu,
    \item úhel došlapu jednotlivých kroků,
    \item náklon trupu a
    \item důležité momenty skoku.
\end{itemize}


\subsection{Uložení výstupu}
\label{ssec:ulozeni}

\paragraph{Parametry}

Hodnoty parametrů vypíšu do souboru pro případnou hlubší biomechanickou analýzu. Jedná se o~\texttt{.csv} soubor, jehož formát je následující.

Parametry ukládám po sloupcích. První sloupec obsahuje relativní čas vůči momentu odrazu v~každém snímku videa, následující sloupce obsahují parametry.

První řádek souboru obsahuje název analyzovaného videa, druhý názvy parametrů, počínaje od druhého sloupce. Ve třetím řádku jsou vypsané jednotky, ve kterých jsou hodnoty parametrů určeny.

Parametry, které reprezentuje jedna hodnota, zabírají jen první řádek hodnot.

Parametry, které neodpovídají snímkům, ale obsahují více hodnot (například doba trvání jednotlivých kroků), jsou uloženy postupně od prvního řádku, dokud jejich hodnoty znám.

Parametry, pro které znám hodnotu v~každém snímku, jsou uloženy v~řádcích odpovídajícím příslušným snímkům videa.

Takto vypadá prvních pár řádků souboru s~hodnotami parametrů po analýze videa \texttt{17.MOV} (viz tabulka \ref{fig:videa}):
\begin{code}[fontsize=\footnotesize]
data/17.MOV
Time,Hips height,Hips velocity loss,Left ankle height,Right ankle height,...
s,px,%,px,px,%,°,s,°,°
-3.85,75.21,106.04,5.62,11.18,27.79,-7.38,0.13,248.43,-2.30
-3.83,79.86,,11.89,3.87,,-0.25,0.20,,1.84
-3.82,75.98,,6.12,7.39,,-14.09,0.03,,-6.30
-3.80,81.02,,8.38,2.88,,21.95,0.03,,-2.21
-3.78,77.05,,6.95,15.49,,29.16,0.08,,-6.30
-3.77,80.58,,18.09,11.37,,-24.36,0.12,,1.84
-3.75,75.86,,4.51,5.81,,-11.59,0.03,,-6.30
\end{code}


\paragraph{Model}

Po analýze videa se výsledný model uloží do textového souboru. První řádek obsahuje název analyzovaného videa.

Následně se v~souboru opakují části reprezentující detekce jednotlivých snímků. Na prvním řádku této části je číslo snímku videa, následuje pozice levého horního rohu snímku v~souřadnicích videa převedených do 3D - hodnota $x$-ové osy roste směrem doprava. Na dalších řádcích jsou vypsány klouby kostry těla v~souřadnicích snímku převedených do 3D. Tyto klouby jsou uloženy postupně, podle zvoleného modelu reprezentace těla, ten je popsán v~sekci \ref{sec:prostredky}.

Takto vypadá prvních pár řádků popisu modelu, který program extrahoval z~videa \texttt{17.MOV} (viz tabulka \ref{fig:videa}):
\begin{code}[fontsize=\footnotesize]
data/17.MOV
0
-719,0,-421
716.211,0,340.27
716.211,0,368.108
710.632,0,373.676
705.053,0,390.378
699.474,0,390.378
727.368,0,373.676
721.789,0,395.946
705.053,0,395.946
710.632,0,407.081
721.789,0,446.054
732.947,0,473.892
716.211,0,412.649
710.632,0,446.054
693.895,0,479.459
716.211,0,390.378
\end{code}

Po analýze modelu se model do souboru neukládá.


\subsection{Zobrazení výstupu uživateli}

Výstup programu je vhodné uživateli přehledně zobrazit. Rozhodl jsem se pro vlastní prohlížeč snímků videa. Mezi jednotlivými snímky se lze pohybovat dopředu i~dozadu. Ve výchozím režimu se zobrazí video se zakreslenou kostrou, jejíž zakreslení do snímku lze vypnout a~znovu zapnout.

Při zobrazení prohlížeče dojde také k~zobrazení snímků, které reprezentují začátek rozběhu, moment odrazu a~moment kulminace nad laťkou.

Při prohlížení daného snímku vypíše program do konzole hodnoty příslušných parametrů. Parametry jsem rozdělil do kategorií podle části skoku, při které jsou podstatné. Například úhel došlapu je irelevantní při skoku, ale je důležitý při rozběhu. Tedy při snímku rozběhu se vypíše například úhel naklonění trupu v~daném snímku a~doba trvání právě probíhajícího kroku. Parametry související s~odrazem vypíše program při zobrazení snímku, který reprezentuje moment odrazu. Způsob zobrazení výstupu je vidět na obrázku \ref{fig:vystup}.

\begin{figure}[p]\centering
    \subfloat[]{
        \includegraphics[width=\textwidth]{vystup1}
    } \\
    \subfloat[]{
        \includegraphics[width=\textwidth]{vystup2}
    }
    \caption{
        \centering\small
        Způsob zobrazení výstupu uživateli. (a)~-~snímky zachycující začátek rozběhu, odraz a~kulminaci boků, prohlížeč videa a~výpis příslušných hodnot analyzovaných parametrů, (b)~-~grafy znázorňující hodnoty parametrů v~průběhu skoku. Pro zobrazení tohoto výstupu byl použit model získaný z~videa \texttt{1.MOV} (viz tabulka \ref{fig:videa}).
    }
    \label{fig:vystup}
\end{figure}


Po ukončení prohlížeče snímků zobrazím uživateli hodnoty parametrů zanesených do grafů, které se načtou z~vygenerovaného souboru.




\section{Programové vybavení}

\subsection{Počítačové vidění a~zpracování obrazu}

Pro potřeby počítačového vidění a~zpracování obrazu lze využít několik knihoven.

\paragraph{OpenCV}

 OpenCV \citep{OpenCV} je jednou z~nejrozšířenějších open-source knihoven. Mezi její výhody patří všestrannost. Je vhodnou volbou pro zpracování obrazu, videí a~vytváření modelů počítačového vidění a~strojového učení. OpenCV vznikla již v~roce 2000, což je jedním z~důvodů, proč ji využívá a~udržuje mnoho vývojářů. Na internetu lze tedy najít mnoho informací o~naprosté většině funkcionality, kterou tato knihovna disponuje.

OpenCV lze spustit na různých platformách a~je vhodnou volbou i~pro mobilní zařízení, což může být užitečné pro budoucí rozšíření aplikace. Tuto knihovnu lze použít v~programu psaném v jazycích C, C++, Python nebo Octave.

\paragraph{OpenVINO}

OpenVINO je knihovna z~dílny společnosti Intel. Existuje ve dvou verzích, jako open-source a~jako distribuce spravovaná právě společností Intel. Knihovna disponuje solidní zásobou modelů hlubokého učení pro potřeby počítačového vidění a~jejich optimalizací pro procesory společnosti Intel.

Ačkoliv se vývojáři snaží podporovat i~jiné procesory, může být spuštění této knihovny na procesorech s~jinou architekturou probematické. Možnosti knihovny pro zpracování obrazu jsou omezené, takže se často používá v~kombinaci s~OpenCV.

API této knihovny je určené pro jazyky C, C++ a~Python.

\paragraph{VisionWorks}

Společnost NVidia vyvinula knihovnu VisionWorks, která využívá rychlosti grafických karet. Tato knihovna slouží především k~detekci a~trasování objektů. Má výborné výsledky v~oblasti autonomního řízení, což není příliš použitelné v~mém programu.

\paragraph{Vision Workbench}

NASA spravuje vlastní knihovnu Vision Workbench, zabývající se počítačovým viděním a~zpracováním obrazu. Hlavním cílem této knihovny je zpracování obrazu a~počítačové vidění pro vesmírné roboty.



\subsection{Grafická reprezentace výstupu}

Užitečným prvkem při analýze videozáznamu je přehrávač videa. Pro mé potřeby by bylo vhodné, aby se pro daný snímek zobrazovaly hodnoty parametrů, takže přehrávač videa nahradím prohlížečem snímků videa.

Pro analýzu parametrů a~jejich případné porovnání je vhodné jejich hodnoty zanést do grafů.

\paragraph{gnuplot}

Jedním z~nejrozšířenějších programů pro tvorbu grafů je gnuplot \citep{gnuplot}. Program se spouští z~příkazové řádky a~běží na několika platformách. Jedná se o~program, jehož první verze vznikla už v~roce 1986, takže má, podobně jako OpenCV, velké množství uživatelů, kteří s~ním mají bohaté zkušenosti.

Gnuplot zpracovává vstup z~příkazové řádky nebo skriptů, které popisují, jaká data se mají vykreslit. Příkazy nejsou psané v~žádném programovacím jazyce, ale mají vlastní syntaxi.

\paragraph{MATLAB}

Velké množstí vědců používá pro práci s~daty MATLAB \citep{MATLAB}. Jedná se o~programovací jazyk, který se specializuje například na numerické výpočty, práci s~maticemi nebo zanášení dat do grafů.

MATLAB je možné provázat s~jinými programovacími jazyky, mezi něž patří C, Java nebo Python.

\paragraph{Matplotlib}

Stále větší oblibě se těší programovací jazyk Python, který disponuje knihovnou pro generování grafů. Tato knihovna se jmenuje Matplotlib \citep{Matplotlib}, která se používá pouze pomocí jazyka Python. Standardním způsobem práce s~touto knihovnou je přes modul Pyplot, jehož funcionalita je podobná MATLABu. Pyplot disponuje množinou funkcí, které umí vykreslit různé typy grafů. Výhodou použití Pythonu je také snadné spouštění skriptů, které není potřeba manuálně kompilovat před jejich spuštěním.

Skripty psané v~Pythonu lze spouštět přímo z~kódu jazyka C++, což umožňuje automatické zobrazení grafů po doběhnutí analýzy videa skoku o~tyči.


































\chapter{Realizace}

\section{Související práce}

\section{Rešerše dostupných knihoven}

\subsection{Počítačové vidění a~zpracování obrazu}

Pro potřeby počítačového vidění a~zpracování obrazu lze využít několik knihoven. Jednou z~nejrozšířenějších open-source knihoven je OpenCV. Mezi její výhody patří všestrannost. Je vhodnou volbou pro zpracování obrazu, videí a~vytváření modelů počítačového vidění a~strojového učení. OpenCV vznikla již v~roce 2000, což je jedním z~důvodů, proč ji využívá a~udržuje mnoho vývojářů. Na internetu lze tedy najít mnoho informací o~naprosté většině funkcionality, kterou tato knihovna disponuje.

OpenCV lze spustit na různých platformách a~je vhodnou volbou i~pro mobilní zařízení, což může být užitečné pro budoucí rozšíření aplikace. Tuto knihovnu lze použít v~programu psaném v jazycích C, C++, Python nebo Octave.

Další možností je využití knihovny OpenVINO. Tato knihovna existuje ve dvou verzích, jako open-source a~distribuce spravovaná společností Intel. Knihovna disponuje solidní zásobou modelů hlubokého učení pro potřeby počítačového vidění a~jejich optimalizací pro procesory společnosti Intel. Ačkoliv se vývojáři snaží podporovat i~jiné procesory, může být spuštění této knihovny na procesorech s~jinou architekturou značně probematické. Možnosti knihovny pro zpracování obrazu jsou omezené, takže se často používá v~kombinaci s~OpenCV.

VisionWorks je zajímavá knihovna od společnosti NVidia. Jedná se o~knihovnu, která využívá rychlosti grafických karet. Tato knihovna slouží především k~detekci a~trasování objektů. Tato knihovna má výborné výsledky v~oblasti autonomního řízení, ale pro mé potřeby není lepším řešením než OpenCV.

NASA spravuje vlastní knihovnu zabývající se počítačovým viděním a~zpracováním obrazu. Hlavním cílem této knihovny je zpracování obrazu a~počítačové vidění pro vesmírné roboty, což není příliš vhodné pro můj projekt.

\subsection{Generování grafů}

Jedním z~nejrozšířenějších programů pro tvorbu grafů je gnuplot. Program se spouští z~příkazové řádky a~běží na několika platformách. Jedná se o~program, jehož první verze vznikla už v~roce 1986, takže má, podobně jako OpenCV, velké množství uživatelů, kteří s~ním mají bohaté zkušenosti.

Gnuplot zpracovává vstup z~příkazové řádky nebo skriptů, které popisují, jaká data se mají vykreslit. Příkazy nejsou psané v~žádném programovacím jazyce, ale mají vlastní syntaxi.

Velké množstí vědců používá pro práci s~daty MATLAB. Jedná se o~programovací jazyk, který se specializuje například na numerické výpočty, práci s~maticemi nebo zanášení dat do grafů. MATLAB je možné provázat s~jinými programovacími jazyky, mezi něž patří C, Java nebo Python.

Stále větší oblibě se těší programovací jazyk Python, který disponuje knihovnou pro generování grafů. Tato knihovna se jmenuje Matplotlib, která se používá pouze pomocí jazyka Python. Standordním způsobem práce s~touto knihovnou je přes modul Pyplot, jehož funcionalita je podobná MATLABu. Pyplot disponuje množinou funkcí, které umí vykreslit různé typy grafů. Výhodou použití Pythonu je také snadné spouštění skriptů, které není potřeba manuálně kompilovat před jejich spuštěním.

Skripty psané v~Pythonu lze spouštět přímo z~kódu jazyka C++, což umožňuje automatické zobrazení grafů po doběhnutí analýzy videa skoku o~tyči.

\section{Návrh řešení}

K~implementaci zpracování obrazu a~počítačového vidění využiji knihovnu OpenCV, program budu psát v~jazyce C++, jelikož je běh výsledného programu rychlejší než při použití Pythonu. Také mám s~jazykem C++ více zkušeností. Pro trasování objektů ve videu použiji \texttt{cv::TrackerCSRT}, který využívá Discriminative Correlation Filter with Channel and Spatial Reliability.

Výstupní soubory popisující model a~parametry skoku zpracuji skripty napsanými v~Pythonu. Výsledné grafy zanesu do grafů s~využitím knihovny Matplotlib.

\subsection{Zpracování videa}

Ke zpracování videa využiji třídu \texttt{cv::VideoCapture}, která otevře video a~uloží jednotlivé snímky videa, aby video nebylo otevřené po celou dobu běhu programu. Snímky videa budou reprezentovat instance \texttt{cv::Mat}, což jsou $n$-dimenzionální pole.

\subsection{Nalezení atletova těla ve videu}

Prvním krokem k~analýze videozáznamu je nalezení pozice atletova těla ve videu. Pro její určení stačí znát pozici ohraničujícího rámečku atletova těla ve snímku, v~němž je poprvé vidět. Může se stát, že v~prvním snímku videa atlet vidět není, a tedy je nutné specifikovat i~číslo snímku, ve kterém atletův ohraničující rámeček znám.

K~zisku těchto informací lze přistoupit dvěma metodami. První z~nich je manuální zadání ohraničujícího rámečku uživatelem. Pro jednoduchst práce s~programem je vhodné použít grafické rozhraní, které dovolí uživateli definovat rámeček pohybem myši po snímku videa. K~tomuto účelu se výborně hodí metoda \texttt{cv::selectROI}.

Druhou metodou nalezení atletovy pozice ve videu je automatická detekce. K~její realizaci nejprve detekuji ohraničující rámečky postav ve snímku videa a~následně budu sledovat jejich pohyb.

Pro detekci postav ve snímku lze použít strukturu \texttt{cv::HOGDescriptor} a~její metodu \texttt{detectMultiScale}, která je schopná detekovat několik postav různých velikostí v~daném snímku. K~detekci postav ve snímku využívá HOG a~SVM.

Detekované postavy je potřeba trasovat i v~následujících snímcích. K~tomu využiji třídu \texttt{cv::Tracker}.

Jelikož je mým cílem detekce atleta, stačí sledovat jen postavy, které se hýbou. Pohyb postavy určím podle vzájemného pohybu ohraničujícího rámečku postavy a~pozadí snímku. Pohyb pozadí ve videu budu trasovat stejně jako pohyb postav. Vyberu si část pozadí a~její pohyb získám s~použitím třídy \texttt{cv::Tracker}.

Atleta odliším od ostatních pohybujících se postav tak, že provede odraz, tedy se směr jeho pohybu změní směrem vzhůru. Jakmile k~odrazu dojde, zapamatuji si číslo snímku, v~němž jsem atleta poprvé detekoval, a~pozici jeho ohraničujícího rámečku v daném snímku.

\subsection{Detekce kostry atleta}

Pro detekci atletovy kostry využiji natrénovaný model projektu OpenPose \citep{OpenPose}. Model je natrénovaný na množině MPII Human Pose Dataset \citep{MPIIHPE}, která obsahuje 25~tisíc obrázků s~více než 40~tisíci postavami s~anotovanými částmi těla. Množina obsahuje postavy vykonávající 410~aktivit, které jsou u~daného obrázku specifikované.

Počet částí těla, které model specifikuje je 16, ale poslední z~nich je pozadí, které budu ignorovat. Tedy se budu zajímat pouze o~prvních 15 částí těla. Jedná se o
\begin{enumerate}
\setcounter{enumi}{-1}
\item hlavu,
\item krk,
\item pravé rameno,
\item pravý loket,
\item pravé zápěstí,
\item levé rameno,
\item levý loket,
\item levé zápěstí,
\item pravou kyčel,
\item pravé koleno,
\item pravý kotník
\item levou kyčel,
\item levé koleno,
\item levý kotník a
\item hrudník.
\end{enumerate}

Model následně načtu do hluboké neuronové sítě \texttt{cv::dnn::Net} pomocí metody \texttt{cv::dnn::readNetFromCaffe}, jelikož je model uložený ve formátu používaného frameworkem Caffe \citep{Caffe}.

Abych mohl síti předat obrázek, je nutné ho nejprve překonvertovat metodou \texttt{cv::dnn::blobFromImage}. Výsledek této konverze je validní vstup pro síť. Jedná se o~čtyřdimenzionální instanci \texttt{cv::Mat}, jejíž rozměry jsou
\begin{itemize}
\item počet vstupních obrázků,
\item počet kanálů vstupních obrázků,
\item výška obrázků a
\item šířka obrázků.
\end{itemize}

Barevné kanály jsou v~rámci knihovny OpenCV i~frameworku Caffe reprezentovány modelem BGR, tedy jsou po složkách uloženy hodnoty modré, zelené a~červené barvy, a~tudíž není nutné řešit konverzi barevných kanálů.

Výsledná instance \texttt{cv::Mat} se následně zpracuje sítí a~vydá výstup, což je opět čtyřdimenzionální instance \texttt{cv::Mat}, jejíž rozměry jsou
\begin{itemize}
\item počet vstupních obrázků,
\item počet výstupních parametrů,
\item výška výstupů a
\item šířka výstupů.
\end{itemize}
Počet výstupních parametrů je 44, ale budu používat pouze prvních 15, které reprezentují části atletova těla potřebné pro vytvoření kostry. Výslednou hodnotou ve výstupu na dané pozici $(i,n,x,y)$ je pravděpodobnost, že bod na řádku $x$, ve sloupci $y$ ve výstupu odpovídá části těla $n$ ve vstupním obrázku $i$. Abych dostal pozici výstupního bodu ve vstupním obrázku, je nutné výstupní souřadnice škálovat poměrem velikostí vstupu a~výstupu sítě.

Jelikož se atlet při skoku otáčí a~model je natrénovaný na postavách, které otočené nejsou, bude pro lepší přesnost detekce potřeba otáčet také vstupní obrázek. Detekci tedy může být vhodné spouštět na více obrázcích zároveň, což síť umožňuje. Vstup sítě vytvořím z~více obrázků metodou \texttt{cv::dnn::blobFromImages}, která k~tomuto účelu slouží.

\subsection{Zisk modelu atleta z~detekované kostry}

\subsubsection{Souřadný systém}

Pozice částí těla atleta budu reprezentovat v~trojrozměrném souřadném systému, jehož počátek bude v~místě kotníku odrazové nohy při odrazu. Ideálním počátkem souřadné soustavy by byl střed zadní hrany kastlíku, ten je bohužel ve videu skrytý za doskočištěm, a~tedy bych jeho pozici musel aproximovat s~pomocí detekce tyče.

Konstantní polohu bude zaujímat pozadí videozáznamu, které bude trasováno instancí třídy \texttt{cv::Tracker}.

Body budu reprezentovat jako uspořádané trojice $(x,y,z)$, kde $x$ bude reprezentovat pozici podle horizontální osy rovnoběžné s~rozběžištěm, $y$ pozici podle horizontální osy kolmé na rozběžiště a~$z$ pozici podle vertikální osy. Hodnoty složek bodů budou růst ve směru rozběhu, od kamery a~vzhůru.

Za jednotku souřadného systému zvolím pixel. Výhodou tohoto řešení je jednodušší převod detekované kostry do souřadného systému a~poměrně přesné porovnání podobných videozáznamů. Nevýhodou je naopak horší přesnost porovnání videí s~různým rozlišením a~jinou vzdáleností kamery od rozběžiště. Tento systém také neřeší vzdálenost atleta od kamery, a~tedy se pohyb atleta na konci rozběhu bude zdát výraznější než na jeho začátku. Tato skutečnost je ještě umocněna faktem, že se atlet nehýbe kolmo k~ose kamery, ale pod jiným úhlem, který horizontální pohyb atleta ještě více potlačuje.

\subsubsection{Přenesení kostry do souřadného systému}

Ve snímku videa je pozice pixelu určena uspořádanou dvojicí $(x,y)$, počátek je v~levém horním rohu a~hodnoty rostou doprava dolů. První složka reprezentuje vzdálenost bodu od levé hrany snímku, druhá vzdálenost bodu od horní hrany snímku.

Při přenesení kostry do souřadného systému bude tedy třeba převrátit hodnoty ve druhé složce. Podle směru atletova rozběhu, který určím podle vzájemného pohybu pozadí videa a~atletova těla, budu muset určit, zda převracet také hodnoty v~první složce. Pokud atlet běží doleva, tak je nutné hodnoty převrátit, pokud poběží doleva, tak nikoliv.

Tímto procesem získám posunuté hodnoty dvou složek ve výsledném souřadném systému. Hodnotu ve složce určené horizontální osou kolmou na rozběžiště nastavím pro jednoduchost na nulu. Hodnotu by bylo možné aproximovat při rozběhu, ale při skoku se atlet otáčí kolem své osy, což zvyšuje obtížnost případné aproximace této hodnoty. Skok může být křivý, tedy atlet nedopadne na doskočiště do osy rozběhu, což je další aspekt, který by bylo nutné vzít v~potaz.

Nakonec bude tedy potřeba posunout hodnoty bodů tak, aby se počátek soustavy posunul do místa kotníku odrazové nohy v~momentu odrazu.

\subsubsection{Určení pohybu jednotlivých snímků}

Kamera není při natáčení skoku o~tyči statická, tedy je nutné vzít v~potaz její pohyb při analýze pohybu atleta. Pro převod kostry detekované ve snímku do souřadného systému bude potřeba zjistit pohyb kamery.

Dostatečným řešením bude určení vzájemné pozice levého horního rohu snímku a pozadí videa. Při zpracování snímku tedy budu znát pozici pozadí v~předešlém a~právě zpracovávaném snímku. Bude stači pozice středu části pozadí, kterou budu sledovat trackerem. Na základě toho, že se pozadí v~souřadném systému nepohybuje, určím vzájemný pohyb levých horních rohů po sobě jdoucích snímků, čímž získám i~vzájemnou polohu atletova těla v~po sobě jdoucích snímcích.

\subsection{Výstup analýzy}

Výsledky analýzy bude potřeba uložit a~zobrazit uživateli. Výsledkem je model a~hodnoty detekovaných parametrů, které uložím do souborů. Soubor popisující model bude možné předat programu jako vstup, aby nemusel v~již analyzovaném videu detekovat atleta znovu, a~provedlo se jen vyhodnocení parametrů.

Jelikož se jedná o~analýzu videa, je vhodné, aby si uživatel mohl prohlédnout video a v~něm detekované body. Tím získá zpětnou vazbu, zda analýza proběhla v~pořádku a~má smysl zkoumat hodnoty parametrů.

Video zobrazím jako prohlížeč jednotlivých snímků videa, k~aktuálně zobrazenému snímku vypíšu hodnoty parametrů, které s~daným snímkem souvisí. Uživatel bude moci zobrazit nebo~skrýt detekce kostry těla.

Po ukončení prohlížeče se zobrazí grafy popisující hodnoty parametrů skoku.

\subsection{Analýza parametrů skoku}



\section{Popis konečného řešení}

\subsection{Struktura programu}














\chapter{Experimenty}

\section{Sběr dat}

Funkčnost programu jsem průběžně ladil na několika videích z~mého archivu. Pro potřeby korektního otestování funkčnosti a~specifických vlastností programu jsem natočil další videa.

Testovací videa obsahují $46$\,\rm videí, které zachycují závodní a~tréninkové skoky a~tréninkové nácviky skoků, videa jsou natočena z různých úhlů, stran a~vzdáleností. Na videích je zachyceno $5$\,\rm atletů ($3$\,\rm muži a $2$\,\rm ženy). K~natočení videí byly využity mobilní telefony Apple iPhone SE ($2020$) a~Samsung Galaxy S10+. Většina videí je natočena v~rozlišení Full HD ($1920\times 1080\,\rm px$), jedno z~videí je po komprimaci, která nastane při poslání videa přes aplikaci WhatsApp. Videa natočená telefonem Apple iPhone SE ($2020$) mají frekvenci snímků $60$\,\rm fps, frekvence snímků videí natočených Samsungem Galaxy S10+ je $30$\,\rm fps. Vliv frekvence snímků videa při analýze je popsán v~sekci \ref{sec:frekvence}.

Seznam videí s~příslušnými vlastnostmi je popsán v~tabulce \ref{fig:videa}.

\begin{figure}[p]\centering
\includegraphics[height=.8\textheight]{videa}
\caption{Testovací videa.}
\small
    A~-~atlet 
        (M$i$~:~muž $i$,
        Ž$i$~:~žena $i$), 
    T~-~typ skoku 
        (0~:~rozběh, 
        1~:~průchod, 
        2~:~průchod na rovné tyči, 
        3~:~zvrat, 
        4~:~zvrat na rovné tyči, 
        5~:~skok, 
        6~:~náskok do ohybu, 
        7~:~proběhnutí, 
        8~:~závodní skok, 
        9~:~závodní skok rozcvičení), 
    V~-~výška laťky~[cm], 
    DT~-~délka tyče~[cm], 
    VÚ~-~výška úchopu~[cm], 
    O~-~vzdálenost místa odrazu od zadní hrany zasouvací skříňky~[cm], 
    DR~-~délka rozběhu od zadní hrany zasouvací skříňky~[m], 
    KR~-~počet kroků, 
    KV~-~počet kroků na videu, 
    D1~-~vzdálenost kamery od osy rozběžiště~[m], 
    D2~-~vzdálenost kamery od zadní hrany zasouvací skříňky~[cm], 
    VK~-~výška kamery nad zemí~[cm], 
    S~-~strana, ze které je video natočeno 
        (P~:~pravá, 
        L~:~levá), 
    R~-~rozlišení videa~[px$\times$px], 
    FPS~-~frekvence snímků videa, 
    K~-~kamera 
        (1~:~Apple iPhone SE (2020),
        2~:~Samsung Galaxy S10+)
\label{fig:videa}
\end{figure}


Hlavní část testovacích dat jsem pořídil na soustředění na ostrově Tenerife. Počasí bylo po celou dobu pobytu příznivé pro skok o~tyči, jelikož svítilo slunce. Z~hlediska testovacích dat to byla škoda, protože tato část testovacích dat neobsahuje videa v~různých světelných podmínkách.

Jedná se o~$32$\,\rm videí, která zachycují skoky tří atletů a~dvou atletek. Kromě samotných skoků jsou na videích zachyceny i~jiné tréninkové prostředky, mezi které patří rozběh s~naznačením odrazu (např. video \texttt{2.MOV}), průchod (pouhý odraz na tyč, tělo zůstává ve víceméně vzpřímené pozici) (např. video \texttt{1.MOV}) nebo náskok do ohybu (průchod, kdy se atlet vrací zpět na rozběh s~pomocí trenéra) (např. video \texttt{3.MOV}). Skoky jsem natáčel z~různých stran, úhlů a~vzdáleností.

Jak jsem popisoval v~sekci \ref{sec:vstup}, závodní skoky jsou typicky natáčeny z~mnohem větší vzdálenosti. Mobilní telefony nedisponují kvalitním přiblížením při natáčení videa, tudíž je atlet na videu mnohem menší. Navíc se na závodech pohybuje více atletů, než na tréninku, i z~tohoto důvodu jsem mezi testovací videa zahrnul mé závodní skoky - $1$\,\rm soutěžní skok z~Mistrovství České republiky mužů a~žen v~hale 2021 a $4$\,\rm rozcvičovací a $5$\,\rm soutěžních skoků ze Zlaté tretry 2021.

Mezi testovací videa jsem dále zahrnul $4$\,\rm tréninkové skoky z~atletické haly na pražském Strahově, mezi kterými jsou $2$\,\rm skoky, které zachycují zvrat na rovné tyči (otočení atleta vzhůru nohama bez ohnutí tyče) (např. video \texttt{33.MOV}).

Úvodní pozice atleta v~testovacích videích, které jsem pro experimenty používal, jsem volil manuálně. Jejich popis je uveden v~tabulce \ref{fig:uvodniramecky}.



\section{Hledání atleta ve videu}

Při spuštění automatického hledání atleta ve všech testovacích videích se atletovo tělo nepodařilo nalézt ve dvou videích (\texttt{5.MOV} a~\texttt{40.mp4}). V~jednom videu označil program za atleta osvětlení stadionu (video \texttt{11.MOV}) kvůli chybné analýze pohybu pozadí, které trasovalo postavu.

Velikosti nalezených atletů se lišily. V~některých videích našel program jen část atletova těla - úvodní rámeček se nacházel jen na trupu, případně nohách atleta. V~některých videích se atletovo tělo detekovalo větší, než bylo třeba. Vliv velikosti úvodních rámečků rozebírám v~sekci \ref{ssec:velikost}.

Atletovo tělo program po nalezení ztratil jen v~nepatrném množství případů. Dokonce ve videích, ve kterých kvůli nevhodnému pozadí selhávalo trasování atleta vybraného manuální cestou. Jelikož se pohyb nalezeného atleta trasuje, je značná šance, že při inicializaci postavy v~prvním snímku dojde k~chybě trasování a~začnou se hledat postavy nové. Problémem tohoto přístupu může být to, že atlet nebude trasován po celou dobu rozběhu, ale jen v~jeho konci.



\section{Volba úvodní pozice atleta}

\subsection{Pozice ohraničujícího rámečku}

Pozice úvodního ohraničujícího rámečku nemá na detekci příliš velký vliv. Část mých testovacích dat je specifická tím, že na druhé straně rozběžiště stojí trenér a~při posunutí úvodního rámečku atleta směrem dozadu se okno pro detekci kostry posouvá více za atleta. Po odrazu, kdy se atletovy nohy pohybují dopředu, obsahuje okno z~větší části trenéra, čímž dochází k~nežádoucí detekci jeho těla. U~ostatních posunutí k~těmto nežádoucím detekcím dochází méně.

Rámečky na těle atleta většinou neobsahují celé tělo, jelikož se atlet přibližuje a~jejich velikost se nemění. Nejčastěji jsou umístěny na atletově trupu a~při zvratu je jejich část pod trupem. Tento problém je možné vyřešit mírným posunutím úvodního rámečku směrem nad nebo před atletovo tělo. Tím se při zvratu atleta rámečky inicializují na části těla, která se pohybuje dopředu a~nahoru. Toto posunuté může ovšem zvýšit míru nežádoucích detekcí, které jsem popsal v~předchozím odstavci.

Pohyb některých částí těla při posunutí úvodního rámečku je vidět na obrázku \ref{fig:posun}.



\subsection{Velikost ohraničujícího rámečku}
\label{ssec:velikost}

Velikost úvodního ohraničujícího rámečku má na detekci značný vliv. Pokud uživatel zvolí příliš malý rámeček, nemusí se do okna určeného pro detekci kostry vejít celé tělo atleta. Aktualizace velikosti detekčního okna probíhá pouze v~případě, že se detekuje celé tělo atleta. Snažím se tak omezit zbytečné zvětšování okna v~případech, kdy se detekuje více postav, které nejsou v~okně celé a~atlet není vidět příliš dobře.

Atlet se v~průběhu videa typicky přibližuje, čímž se jeho pozice ve videu zvětšuje, ale velikost mřížky zůstává stejná jako na začátku. To může zapříčinit jen částečné detekce atletova těla a~tedy se okno pro detekci nikdy nezvětší.

Pokud uživatel zvolí naopak větší rámeček, může se stát, že se trackery mřížky chytí ostatních objektů. Pokud je video pořízené dostatečný čas před začátkem atletova rozběhu, mohou se trackery pohybovat s~postavami v~pozadí a~ztratit tak atleta hned na začátku.

Příliš velký úvodní rámeček může způsobit, že detektor kostry detekuje kostru jiné postavy nebo smíchá klouby různých postav v~jednu kostru. [OBRÁZEK] Tyto detekce neovlivní chování trackeru, ale značně znepřesňují analýzu skoků. Toto chování blíže popisuji v~sekci \ref{sec:postavy}.




\section{Pohyb kamery}

\subsection{Vertikální pohyb kamery}

Vertikální pohyb kamery nemá na výsledné hodnoty parametrů téměř žádný vliv.

Pohyb atleta ve videu určuji podle pozice detekované části pozadí ve videu, tudíž při se při pohybu kamerou ve snímku pohybuje pozadí opačným směrem. Tato vlastnost stabilizuje souřadný systém videa při pohybech kamery. Stabilizace je vidět v~grafech na obrázku \ref{fig:vpohyb}. V~grafech je znázorněna pozice levého a~pravého kotníku ve videích \texttt{17.MOV} a~\texttt{18.MOV}. Kamera se při pořizování videa \texttt{17.MOV} pohybovala v~průběhu první poloviny rozběhu směrem vzhůru, v~průběhu druhé poloviny rozběhu směrem dolů. Záběr tedy tvořil oblouk, jelikož zároveň trasoval atleta. Při pořizování druhého videa k~vertikálním pohybům kamery nedocházelo.

\begin{figure}[h]\centering
    \subfloat{
        \includegraphics[width=.45\textwidth]{vpohybl}
    }
    \subfloat{
        \includegraphics[width=.45\textwidth]{vpohybp}
    }
    \caption{Vliv horizontálního pohybu kamery na souřadný systém.}
    \small
    Grafy znázorňují pozici levého a~pravého kotníku ve videích \texttt{17.MOV} a~\texttt{18.MOV}.
    \label{fig:vpohyb}
\end{figure}

Pokud bych souřadný systém nestabilizoval trasováním pozadí, vytvořil by graf pozic částí atletova těla při analýze videa \texttt{17.MOV} oblouk opačným směrem, než kterým se pohybovala kamera. Takovýto oblouk se v~grafu nenachází a~výsledné parametry nejsou nijak ovlivněny.



\subsection{Rotace kamery}

Zajímavým experimentem je analýza videí, při jejichž natáčení jsem rotoval s~kamerou o~$90\,\rm \degree$ po (\texttt{21.MOV}) a~proti (\texttt{22.MOV}) směru hodinových ručiček. Videa jsou pořízena z~levé strany rozběžiště.

Na trasování atleta ani detekci kostry rotace vliv nemá. Je to dáno také charakteristickou vlastností skoku o~tyči a~to rotací atletova těla. S~tou je nutné při analýze videozáznamu skoku o~tyči počítat.

Rotace kamery znamená také rotaci souřadného systému videa, na kterou není připraven. Z~pohybu horizontálního se tedy postupně stává pohyb vertikální. Tuto rotaci souřadného systému zobrazuje graf na obrázku \ref{fig:rotace}. Detekce částí těla proběhla správně na téměř všech snímcích videa. Problematická byla jen při dopadu do doskočiště, kdy byl atlet schovaný za stojany.

\begin{figure}[h]\centering
    \includegraphics[width=\textwidth]{rotace}
    \caption{Vliv rotace na pozici modelu v~prostoru.}
    \small
    Graf znázorňuje pozici hrudníku při analýze videí \texttt{21.MOV} a~\texttt{22.MOV}.
    \label{fig:rotace}
\end{figure}

Jelikož se v~obou videích trasovala část pozadí, která byla před atletem (na snímcích videa se tedy nacházela vlevo), měnila se s~rotací také pozice atleta ve vertikálním směru. Při rotaci po směru hodinových ručiček (video \texttt{21.MOV}) posouvala rotace atleta pod trasovanou část pozadí, což mělo za důsledek klesání atleta v~souřadném systému. Odraz se projevil horizontálním pohybem v~původním směru rozběhu. Při opačné rotaci (video \texttt{22.MOV}) se atlet začal postupně pohybovat nad trasované pozadí, tedy začal stoupat a~odraz se projevil horizontálním pohybem proti původnímu směru rozběhu. Míra vertikálního posunu atleta při rotaci je dána vzájemnou vzdáleností atletova těla a~trasované části pozadí.

Parametry získané analýzou těchto videí tedy nemají vysokou vypovídající hodnotu.




\section{Pozice kamery}

\subsection{Směr rozběhu}

Právě kvůli porovnání skoku z~různých stran jsem pořídil videa \texttt{27\_L.mp4} a~\texttt{27\_P.MOV}. Jedná se o~tentýž skok, který je natočený ze stejné vzdálenosti od rozběžiště a~velice podobné vzdálenosti od zadní hrany zasouvací skříňky. Videa mají rozdílnou frekvenci snímků, ale stejné rozlišení, tudíž je atlet na obou videích v~odpovídajících fázích skoku stejně veliký.

Do grafů na obrázku \ref{fig:strany} jsem zakreslil pozice levé a~pravé kyčle. Pro lepší porovnání jsem posunul souřadné systémy podle přesného momentu odrazu. Ten nastává pro video \texttt{27\_L.mp4} na snímku číslo $101$ a~ve videu \texttt{27\_P.MOV} na snímku číslo $284$. Program určil odraz ve snímcích číslo $109$ a~$1$ (počítáno od $0$). Chybná detekce místa odrazu byla způsobena nepřesnou implementací filtrování kroků, vylepšení chybovosti navrhuji v~sekci \ref{ssec:vylepseni}.

\begin{figure}[h]\centering
    \subfloat{
        \includegraphics[width=.45\textwidth]{27_lkycel}
    }
    \subfloat{
        \includegraphics[width=.45\textwidth]{27_pkycel}
    }
    \caption{Vliv horizontálního pohybu kamery na souřadný systém.}
    \small
    Grafy znázorňují pozici levé a~pravé kyčle ve videích \texttt{27\_L.mp4} a~\texttt{27\_P.MOV}.
    \label{fig:strany}
\end{figure}

Z~grafů je vidět, že jsou pozice vyznačených částí těla velice podobné.

Ve fázi rozběhu jednotlivých kroků se kyčle pohybovaly pod jiným úhlem. Vliv na tento úhel má například mírná rotace kamery, jelikož nebyly kamery na stativu, nelze vliv rotace ignorovat.

Z~grafů je patrné, že detekce rozběhu proběhla lépe na videu \texttt{27\_P.MOV}, ale tělo ve fázi skoku se podařilo lépe detekovat na videu \texttt{27\_L.mp4}.

Tělo se při rozběhu na videu natočeném z~pravé strany podařilo detekovat lépe díky lepšímu kontrastu s~pozadím, vliv kontrastu probírám více v~sekci \ref{sec:kontrast}. Při fázi skoku byla pro detekci výhodnější levá strana, jelikož je z~této strany lépe vidět levá paže (která je z~pravé strany schovaná za tělem) a~při obratu, který atlet na tyči provádí, je atlet natočen čelem k~levé kameře.



\subsection{Úhel záběru}

[Porovnání skoku natočeného z jiných úhlů]



\section{Další postavy ve videu}
\label{sec:postavy}

Další postavy ve videu mají na detekci vliv, zvlášť když se nachází blízko atleta. Tato problematika je vidět na obrázku \ref{fig:spatnakostra}.

Výsledné hodnoty parametrů detekce špatné kostry značně znehodnocuje.




\section{Splynutí s~pozadím}
\label{sec:kontrast}

Pro trasování atleta a~detekci kostry je vhodné, aby byl atlet od pozadí dobře rozpoznatelný. Pokud atlet s~pozadím splývá, je pro tracker i~detektor částí těla složité správně určit jeho pozici. Tento problém je výborně vidět na obrázku \ref{fig:kontrast}.

\begin{figure}[p]\centering
    \subfloat{
        \includegraphics[width=.45\textwidth]{kontrast_b1}
    }
    \subfloat{
        \includegraphics[width=.45\textwidth]{kontrast_o1}
    } \\
    \subfloat{
        \includegraphics[width=.45\textwidth]{kontrast_b2}
    }
    \subfloat{
        \includegraphics[width=.45\textwidth]{kontrast_o2}
    } \\
    \subfloat{
        \includegraphics[width=.45\textwidth]{kontrast_b3}
    }
    \subfloat{
        \includegraphics[width=.45\textwidth]{kontrast_o3}
    } \\
    \small
    Snímky znázorňují ztrátu atleta při jeho splynutí s~pozadím (vlevo) a~úspěšné trasování, pokud je atlet od pozadí dobře rozeznatelný (videa \texttt{36.mp4} a~\texttt{40.mp4}. 
    \label{fig:kontrast}
\end{figure}

V~případě trackeru se jeho ohraničující rámeček v~některých případech chytí na pozadí videa a~tedy ztratí atleta.

Detekce kostry má větší problém s~určením správné pozice kloubů těla, pokud je pozadí blízko atleta členité. Na obrázku \ref{fig:vlivpozadi} [doplním] je vidět několik příkladů tohoto chování detektoru.




\section{Frekvence snímků}
\label{sec:frekvence}

Frekvence snímků videa (pokud se pohybuje v~rozumných hodnotách) má na výsledek detekce zanedbatelný vliv. Značný vliv může mít na přesnost výpočtu hodnot biomechanických parametrů skoku, jelikož na ní do jisté míry závisí.

Z~hlediska časové náročnosti analýzy videa je zajímavé sledovat vliv frekvence snímků videa na přesnost detekcí. Vynecháním poloviny snímků videa trvá analýza přibližně polovinu času, což může být výhodné pro analýzu videí, která obsahují spoustu snímků - ať už z~důvodu vysoké frekvence či délky videa.

Jelikož jsou mezi testovacími daty i~videa s~frekvencí snímků $30$\,\rm fps, nepředpokládal jsem, že trasování selže.

Problematičtější mohla být detekce kostry. S~některými fázemi skoku má detektor značné problémy, ty popisuji v~sekci \ref{sec:shrnuti}. Některé snímky se detektoru podaří detekovat, ale úspěšnost není vysoká. Může se tedy stát, že by korektní detekce proběhly ve vynechaných snímcích. Detekce kostry spoléhá na natočení okna, které určuje detekce předchozí kostry, což může být z~hlediska vynechání snímků, v~nichž by k~detekci došlo, velice problematické. Okno pro detekci se v~takovém případě neotočí, a~tedy nedetekuje postavu atleta správně.

Video, které jsem zvolil pro porovnání, je \texttt{23.MOV}. Toto video je pořízeno frekvencí $60$\,\rm fps, analýzu jsem spustil bez úprav a~poté jsem z~videa vynechal každý druhý snímek. Nižší frekvenci jsem netestoval, jelikož $30$\,\rm snímků za vteřinu je na spodní hranici přípustné frekvence pro analýzu videozáznamu skoku o~tyči, jelikož jsou pohyby při skoku velice rychlé.

Pozici atleta ve videu jsem vybral ručně v~prvním snímku, který byl zahrnutý v~obou videích, tedy inicializace atleta proběhla v~obou videích stejně.

Jak je vidět v~grafu naklonění trupu získaného analýzou videa \texttt{30.mp4} \ref{fig:30_naklon}, tento problém není příliš podstatný a~vliv na detekci má spíše póza, kterou atletovo tělo zaujímá, než jeho samotná rotace.

\begin{figure}[h]\centering
    \includegraphics[width=\textwidth]{30_naklon}
    \caption{Náklon trupu.}
    \small
    Při analýze videa \texttt{30.mp4} dokázal program správně detekovat atleta, jehož trup byl oproti vzpřímené pozici natočen o~$-150\,\rm \degree$. Program v~té chvíli předpokládal poslední známý náklon trupu, tedy přibližně $-20\,\rm \degree$ (záporná hodnota značí záklon).
    \label{fig:30_naklon}
\end{figure}

Výsledky experimentu s~vynecháním snímků jsou zobrazené na obrázku \ref{fig:fps}. Na grafech je vidět, že detekce selhávala ve velice podobných místech a~frekvence videa na výslednou detekci značný vliv neměla. Také nepřesnosti detekcí, které jsou vidět na grafu pohybu pravého zápěstí, jsou velice podobné.

\begin{figure}[h]\centering
    \subfloat{
        \includegraphics[width=.45\textwidth]{23_fps_hlava}
    }
    \subfloat{
        \includegraphics[width=.45\textwidth]{23_fps_pzapesti}
    } \\
    \subfloat{
        \includegraphics[width=.45\textwidth]{23_fps_lkycel}
    }
    \subfloat{
        \includegraphics[width=.45\textwidth]{23_fps_pkycel}
    }
    \caption{Vliv frekvence snímků na detekci.}
    \small
    Na grafech jsou znázorněny pozice hlavy, pravého zápěstí, levé a~pravé kyčle ve videu \texttt{23.MOV} bez úpravy frekvence snímků a s~poloviční frekvencí snímků (každý druhý snímek je vynechaný).
    \label{fig:fps}
\end{figure}

Jelikož se jednalo o~stejné video se stejnou úvodní pozicí atleta, rozhodl jsem se počátek souřadného systému pro porovnání pozic jednotlivých částí těla nepřesouvat do pozice odrazového kotníku v~momentu odrazu, ale nechal jsem jej ve středu úvodního ohraničujícího rámečku. Posun by mohl být různý, pokud by se detekce lišily a~výsledné grafy by nemusely být tak přehledné.

[PARAMETRY]




\section{Shrnutí}
\label{sec:shrnuti}

Shrnutí analýzy testovacích videí.



\subsection{Navrhovaná vylepšení}

Na závěr bych rád popsal, jakými způsoby by bylo možné funkcionalitu programu vylepšit.


\subsubsection{Trasování těla}

Úspěšnosti trasování těla atleta při skoku by mohlo pomoci její natáčení podle fáze, ve které se skokan nachází. Přesnějším způsobem trasování může být kombinace implementovaného přístupu s~využitím mřížky a~metody otáčení snímku, kterou program používal předtím. Mřížku by bylo možné otáčet podle úhlu naklonění trupu.


\subsubsection{Detekce kostry}

Ideálním přístupem k~detekci pozice těla pro takto specifické využití je natrénování detektoru na vhodné množině dat. K~detekci postavy pouze při skoku o~tyči by většinu takové množiny dat měly tvořit právě snímky zachycující skok o~tyči. Databáze MPII Human Pose dataset \citep{MPIIHPE}, na které je natrénovaný model, který aplikace nyní využívá obsahuje naprosté minimum dat týkajících se skoku o~tyči, aby dobře generalizoval. Pro účely mé aplikace není vysoká míra generalizace nutná, při detekci těla atleta, když se ve snímku poblíž nachází jiná postava přesnosti spíše škodí.

Použití vhodné trénovací databáze by mohlo omezit potřebu určení pozice atleta ve videu pomocí trackerů.

Kromě trénování specializovaného modelu by mohlo být užitečné zkoumat celý výstup sítě, která detekci zajišťuje. Tento způsob by mohl omezit případy, kdy se detekuje více postav najednou (viz sekce \ref{sec:postavy}), ale nepomohla by s~neúspěšnými detekcemi ve fázi skoku.

Při popisu experimentu v~sekci \ref{ssec:velikost} jsem popsal problém s~volbou příliš malého úvodního rámečku atletova těla. Menší rámeček občas zvolí automatické hledání atleta, což může mít za důsledek neúplnou detekci kostry. Tento problém by se dal snadno vyřešit aktualizací velikosti detekčního okna bez ohledu na úspěšnost detekce. Při implementaci jsem dal přednost ošetření případu, který nastává velmi vzácně, na rozdíl od volby příliš malého ohraničujícího rámečku.

Detekce kostry dosahuje horších výsledků z~důvodu zmenšení rozlišení okna použitého pro její detekci. Toto zmenšení má obrovský vliv na rychlost zpracování videa, tudíž jsem k~němu přistoupil. Ponechání původního rozlišení by tedy mělo přesnost detekce mírně zlepšit.


\subsubsection{Časová náročnost}

Rozdělení videa do několika částí a~jejch následné paralelní zpracování by časovou náročnost analýzy značně snížilo. Dalším přístupem by mohlo být rozdělení detekce kostry do více částí a~jejich paralelní zpracování.

Kromě využití paralelismu by mohl pomoci specializovaný model pro detekci částí těla, jelikož nyní při neúspěšné detekci zkouším další natočení snímku, která v~určitých případech nepomohou. Tento model by žádné další natočení nepotřeboval, a~tudíž by se detekce spouštěla na menším množství vstupů.

Časovou náročnost při detekci postav zbytečně zvyšuje množství instancí třídy \texttt{background\_tracker}, díky nimž určuji pohyb postav ve videu. Tyto trackery často trasují stejnou část pozadí, ale aktualizují se pro každou instanci. Původně jsem tracker pozadí implementoval v~závislosti na poloze příslušné postavy a~objektový návrh jsem poté nestihl změnit.


\subsubsection{Analýza parametrů}

Přesnost analýzy parametrů, ale také zanesení modelu atleta do souřadného systému, ovlivňuje rotace kamery. Stabilizaci této rotace by zajistila dvojice trackerů pozadí, jejichž vzdájemná pozice by dokázala určit míru rotace jednotlivých snímků. Díky této hodnotě bych snadno určil přesnou pozici atleta v~prostoru i~při rotaci kamery.


\chapter*{Závěr}
\addcontentsline{toc}{chapter}{Závěr}

Jednotlivé kroky analýzy videozáznamu skoku o~tyči, mezi které patří extrakce modelu atleta, detekce důležitých momentů skoku, vyhodnocení a~zobrazení parametrů se mi podařilo implementovat. Ačkoliv není přesnost analýzy vždy dostatečná, jedná se o~úvodní krok k~realizaci aplikace, která by videozáznamy dokázala spolehlivě a~přesně analyzovat a~na základě výsledků těchto analýz skoky porovnávat, což byl můj prvotní cíl.

Původně jsem měl v~úmyslu vytvořit aplikaci pro mobilní zařízení, která by videa analyzovala v~průběhu tréninků a~soutěží. Převod aplikace z~počítače na platformu iOS byl složitější, než jsem očekával a~raději jsem svou pozornost zaměřil na vylepšení funkcionality počítačové aplikace. Posléze se ukázalo, že časová náročnost analýzy by bez modifikací implementovaného algoritmu nezvládla videa analyzovat včas.

Značnou pozornost jsem při psaní práce věnoval automatickému nalezení atleta ve videu. Parametry pro vyhledávání postav ve videu se mi bohužel nepodařilo nastavit tak, abych byl s~úspěšností vyhledávání spokojen. Proto je manuální nalezení atleta ve videu výchozím nastavením programu.

Způsob trasování atleta jsem v~průběhu psaní práce několikrát měnil, úspěšnost se stále zlepšovala, ale ani aktuální metoda není $100\%$ úspěšná ve fázi skoku.

Detekce částí atleta má problémy s~malými postavami a~při skoku, především ve fázi skoku, kdy má atlet nohy blízko trupu. Tyto nepřesnosti jsou důsledkem zvoleného modelu, který jednotlivé části těla detekuje a~jeho trénovací databází. Ani pokusy o~natočení snímku detekci příliš nepomáhají.

Na přesnostech detekce atletova těla závisí hodnoty zkoumaných parametrů (a~také důležitých momentů skoku). Například detekce momentu odrazu spoléhá na kvalitní detekci atletových nohou, která ne vždy proběhne. Kvalitní analýza parametrů si zakládá na důkladné znalosti biomechanických principů a~její implementace vyžaduje mezioborovou spolupráci.

Věřím, že jsem svou prací ukázal, že je možné vytvořit program schopný kvalitní analýzy parametrů skoku o~tyči, kterou by bylo možné využívat jak v~tréninkovém, tak závodním prostředí.




















































%%% Seznam použité literatury
%%% Seznam použité literatury (bibliografie)
%%%
%%% Pro vytváření bibliografie používáme bibTeX. Ten zpracovává
%%% citace v textu (např. makro \cite{...}) a vyhledává k nim literaturu
%%% v souboru literatura.bib.
%%%
%%% Příkaz \bibliographystyle určuje, jakým stylem budou citovány odkazy
%%% v textu. V závorce je název zvoleného souboru .bst. Styly plainnat
%%% a unsrt jsou standardní součástí latexových distribucí. Styl czplainnat
%%% je dodáván s touto šablonou a bibTeX ho hledá v aktuálním adresáři.

\bibliographystyle{czplainnat}    %% Autor (rok) s českými spojkami
% \bibliographystyle{plainnat}    %% Autor (rok) s anglickými spojkami
% \bibliographystyle{unsrt}       %% [číslo]

\renewcommand{\bibname}{Seznam použité literatury}

%%% Vytvoření seznamu literatury. Pozor, pokud jste necitovali ani jednu
%%% položku, seznam se automaticky vynechá.

\bibliography{literatura}

%%% Kdybyste chtěli bibliografii vytvářet ručně (bez bibTeXu), lze to udělat
%%% následovně. V takovém případě se řiďte normou ISO 690 a zvyklostmi v oboru.

% \begin{thebibliography}{99}
%
% \bibitem{lamport94}
%   {\sc Lamport,} Leslie.
%   \emph{\LaTeX: A Document Preparation System}.
%   2. vydání.
%   Massachusetts: Addison Wesley, 1994.
%   ISBN 0-201-52983-1.
%
% \end{thebibliography}


%%% Obrázky v bakalářské práci
%%% (pokud jich je malé množství, obvykle není třeba seznam uvádět)
\listoffigures

%%% Tabulky v bakalářské práci (opět nemusí být nutné uvádět)
%%% U matematických prací může být lepší přemístit seznam tabulek na začátek práce.
%\listoftables

%%% Použité zkratky v bakalářské práci (opět nemusí být nutné uvádět)
%%% U matematických prací může být lepší přemístit seznam zkratek na začátek práce.
%\chapwithtoc{Seznam použitých zkratek}

%%% Přílohy k bakalářské práci, existují-li. Každá příloha musí být alespoň jednou
%%% odkazována z vlastního textu práce. Přílohy se číslují.
%%%
%%% Do tištěné verze se spíše hodí přílohy, které lze číst a prohlížet (dodatečné
%%% tabulky a grafy, různé textové doplňky, ukázky výstupů z počítačových programů,
%%% apod.). Do elektronické verze se hodí přílohy, které budou spíše používány
%%% v elektronické podobě než čteny (zdrojové kódy programů, datové soubory,
%%% interaktivní grafy apod.). Elektronické přílohy se nahrávají do SISu a lze
%%% je také do práce vložit na CD/DVD. Povolené formáty souborů specifikuje
%%% opatření rektora č. 72/2017.
\appendix
\chapter{Přílohy}

\section{První příloha}

\openright
\end{document}
